    \documentclass[lettersize,journal]{IEEEtran}
    
    % ---- Unicode sanitization (IMPORTANT) ----
    \DeclareUnicodeCharacter{202A}{}
    \DeclareUnicodeCharacter{202C}{}
    \DeclareUnicodeCharacter{200E}{}
    \DeclareUnicodeCharacter{200F}{}
    
    % ---- Math & theorem support ----
    \usepackage{amsmath,amssymb,amsfonts}
    \usepackage{amsthm}
    \usepackage{bm}
    \usepackage{algorithm}
    \usepackage{algpseudocode}
    
    \newtheorem{theorem}{Theorem}
    \newtheorem{lemma}{Lemma}
    \newtheorem{definition}{Definition}
    \newtheorem{game}{Game}
    \usepackage{amssymb}
    
    % ---- Algorithms ----
    \usepackage{algorithm}
    \usepackage{algpseudocode}
    
    % ---- Figures & tables ----
    \usepackage{graphicx}
    \usepackage{subcaption}
    \usepackage{booktabs}
    \usepackage{tabularx}
    \usepackage{multirow}
    \usepackage{array}
    \usepackage{adjustbox}
    
    \usepackage{array}
    \usepackage{pifont}
    \newcommand{\cmark}{\ding{51}} % checkmark
    \newcommand{\xmark}{\ding{55}} % cross mark
    
    % ---- Text & formatting ----
    \usepackage{cite}
    \usepackage{textcomp}
    \usepackage{url}
    \usepackage{verbatim}
    \usepackage{stfloats}
    \usepackage{enumitem}
    \usepackage{setspace}
    \usepackage{bm}
    \usepackage{xcolor}
    \usepackage{listings}
    \usepackage{makecell}
    
    % ---- References ----
    \usepackage[capitalize]{cleveref}
    
    \hyphenation{op-tical net-works semi-conduc-tor IEEE-Xplore}
    
    % updated with editorial comments 8/9/2021
    
    \begin{document}
    
    \title{VeRedact-PQ: Scalable Post-Quantum Redaction with Multi-Party Authorization and Verifiable Auditing for Permissioned Blockchains}
    
    %}
    %
    %\thanks{
    %Corresponding author: 
    %Somchart Fugkeaw
    %(e-mail: [somchart@siit.tu.ac.th]
    
    % The paper headers
    \markboth{Journal of \LaTeX\ Class Files,~Vol.~14, No.~8, August~2021}%
    {Shell \MakeLowercase{\textit{et al.}}: A Sample Article Using IEEEtran.cls for IEEE Journals}
    
    %\IEEEpubid{0000--0000/00\$00.00~\copyright~2021 IEEE}
    % Remember, if you use this you must call \IEEEpubidadjcol in the second
    % column for its text to clear the IEEEpubid mark.
    
    \author{
    \IEEEauthorblockN{Rangsimann Sattayarom, Ratchanon Wongwitutai, Natthawat Tungsriworakan, \\
    Tagrid Chongkolrattanapond, Somchart Fugkeaw\\}
    \IEEEauthorblockA{Sirindhorn International Institute of Technology, Thammasat University, Thailand\\
    6622771747@g.siit.tu.ac.th, 
    6622780268@g.siit.tu.ac.th,
    6622772364@g.siit.tu.ac.th,
    6622781175@g.siit.tu.ac.th,
    somchart@siit.tu.ac.th}
    }
    
    \maketitle
\begin{abstract}

Redactable permissioned blockchains enable controlled modification of
committed transactions, but existing approaches face challenges in
distributed authorization, post-quantum security, verifiable
accountability, and scalability under high redaction workloads.
Processing each redaction independently can repeatedly incur
policy-validation, committee, authenticated-state update, and audit
costs. This paper proposes \emph{VeRedact-PQ}, a scalable post-quantum
framework for policy-bound redaction and verifiable auditing in
permissioned blockchains. VeRedact-PQ combines state-bound
post-quantum zero-knowledge validation with threshold committee
authorization and distributed post-quantum chameleon-hash adaptation.
Its Cost- and Overlap-Aware Redaction Aggregation (CORA) identifies
validated requests that can share execution work, while Shared-Frontier
Merkle Redaction (SFMR) exploits this locality through shared Merkle
processing and batch-level chameleon-hash adaptation. Policy-Bound
Batch Redaction Proofs (PBRP) bind request validation and committee
authorization to the finalized state transition, preserving
request-level accountability under shared execution. For auditing, a
sharded authenticated redaction index and State-Aware Verification
Evidence Reuse (SAVER) organize evidence as state-bound dependencies,
allowing common authorization, transition, and checkpoint evidence to
be verified once under matching contexts while selectively escalating
challenged requests to independent post-quantum zero-knowledge
verification. Experimental evaluation examines end-to-end redaction
performance, multi-party authorization, audit and verification
efficiency, selective verification, and permissioned-blockchain
finalization overhead. VeRedact-PQ thereby reduces repeated
cryptographic work across the redaction lifecycle without relaxing
request-level policy enforcement or state-transition accountability.

\end{abstract}

\begin{IEEEkeywords}
Permissioned Blockchain, Redactable Blockchain, Post-Quantum Security,
Post-Quantum Chameleon Hash, Verifiable Redaction, Privacy-Preserving
Auditing, Scalable Redaction.
\end{IEEEkeywords}

\section{Introduction}

Permissioned blockchains provide a shared and auditable ledger for
multi-organization environments without requiring complete mutual
trust~\cite{ref7}. Their integrity relies on cryptographic immutability:
once a transaction is committed, unauthorized modification becomes
detectable. Strict immutability, however, can conflict with legitimate
requirements such as correcting erroneous records, removing sensitive
information for privacy or regulatory compliance, and updating obsolete
data~\cite{ref9,ref21,ref22,ref38}. The challenge is therefore not
simply how to modify blockchain data, but how to ensure that each
modification is \emph{policy-compliant, jointly authorized,
state-consistent, and auditable} without sacrificing scalability.

Consider a consortium blockchain receiving hundreds of compliance-driven
redaction requests across different transaction batches. Some requesters
may prove eligibility using private credentials. Even after validation,
processing requests independently repeats committee approvals, Merkle
updates, and chameleon-hash adaptations. Naive batching, however, risks
executing stale or unauthorized requests under authorization intended
for valid ones. Later, an auditor examining many redactions may again
verify identical authorization, transition, and checkpoint evidence.
Scalable redaction therefore requires safe sharing of cryptographic work
across \emph{validation, authorization, execution, and verification}
without expanding redaction authority.

Chameleon hashes are widely used to modify committed data while
preserving blockchain linkage~\cite{ref1,ref2,ref6,ref30}. Subsequent
schemes introduce fine-grained access control~\cite{ref8,ref15},
decentralized trapdoor management~\cite{ref3,ref19}, dynamic
updates~\cite{ref5,ref16}, threshold redaction~\cite{ref27}, and
traceable or publicly auditable redaction~\cite{ref33,ref34,ref35}.
Nevertheless, preserving a ledger commitment does not prove that a
modification was authorized for the particular requester, operation,
policy, and current state. Because an authorized chameleon-hash
collision intentionally preserves ledger consistency, conventional
chain verification alone cannot establish \emph{why} a redaction was
allowed or whether the authorized request matches the modified state.

Privacy and long-term security introduce further requirements.
Fine-grained and policy-hiding schemes restrict redaction authority or
conceal policies~\cite{ref4,ref18,ref31,ref34}, while verifiable designs
support redaction auditing~\cite{ref13,ref20,ref33}. However,
authorization may depend on private credentials that should not be
disclosed, and auditors should not require sensitive transaction
content. Meanwhile, quantum-resistant chameleon hashes have been
studied for redactable blockchains~\cite{ref10,ref11,ref17,ref37}.
Protecting only the chameleon hash is insufficient if requester
authentication, private-policy verification, committee approval, or
audit authentication remains quantum-vulnerable.

Scalability must therefore be addressed jointly with authorization and
verifiability. Prior work considers scalable redaction
~\cite{ref14,ref21,ref26}, efficient data-structure updates
~\cite{ref23}, lightweight verification~\cite{ref20}, and robust
threshold redaction~\cite{ref27}. Yet request-by-request processing
repeats policy, committee, Merkle, and chameleon-hash operations, while
conventional batching does not exploit whether requests affect the same
transaction batches or share Merkle paths. Similarly, independent
auditing repeatedly verifies evidence common to multiple redactions.
The key problem is therefore to \emph{identify, preserve, and safely
reuse cryptographic sharing} throughout the redaction lifecycle while
retaining request-level authorization and state binding.

To address this problem, we propose \emph{VeRedact-PQ}, a scalable
post-quantum and verifiable redaction framework for permissioned
blockchains. VeRedact-PQ first binds each request to its authenticated
transaction and policy state through PQ authentication and PQZK-based
private-policy validation. It then introduces a coordinated optimization
pipeline: \emph{Cost- and Overlap-Aware Redaction Aggregation (CORA)}
identifies validated requests that can share execution work;
\emph{Shared-Frontier Merkle Redaction (SFMR)} realizes this sharing
during state modification; and \emph{State-Aware Verification Evidence
Reuse (SAVER)} preserves and reuses shared dependencies during auditing.
A \emph{Policy-Bound Batch Redaction Proof (PBRP)} links validation,
committee authorization, and the finalized transition.

The main contributions are summarized as follows:

\begin{itemize}

\item \textbf{State-Bound Post-Quantum Redaction.}
We design an end-to-end security-critical redaction path integrating
PQ requester authentication, state-bound PQZK private-policy
verification, $t$-of-$n$ PQ committee authorization, and distributed
Post-Quantum Chameleon Hash (PQCH) adaptation. The
\emph{Redaction-Aware Dual Commitment (RADC)} separates immutable
transaction-policy context from mutable content state, while SA-RLI
supports authenticated state resolution.

\item \textbf{Cost- and Overlap-Aware Redaction Aggregation.}
We introduce \emph{CORA}, which groups already validated requests using
workload conditions and authenticated execution locality. Unlike
load-only batching, CORA considers affected transaction batches and
shared Merkle frontiers to reduce subsequent cryptographic work, while
bounded waiting preserves fairness without allowing locality to
determine authorization eligibility.

\item \textbf{Shared-Frontier Verifiable Redaction.}
We develop \emph{SFMR} to jointly execute authorized modifications
sharing Merkle authentication structure, processing overlapping paths
once and requiring one distributed PQCH adaptation per affected batch.
PBRP binds request-level validation, threshold authorization, the
authorization-bound SFMR transition, and finalized checkpoint,
preserving individual verifiability under coalesced execution.

\item \textbf{State-Aware Verification Evidence Reuse.}
We introduce \emph{SAVER}, which organizes redaction auditing as
dependency-aware verification over content-addressed request,
authorization, transition, and checkpoint evidence. Shared evidence is
verified once and reused only under an identical authenticated state
context. Challenged requests can selectively escalate to independent
PQZK verification without revalidating unaffected shared evidence,
while RAI provides authenticated historical query resolution.

\end{itemize}


    
\section{Related Work}
   \begin{table*}[t]
\centering
\caption{Comparison of VeRedact-PQ with Representative Redactable
Blockchain Schemes}
\label{tab:comparison}
\renewcommand{\arraystretch}{1.15}
\setlength{\tabcolsep}{4.0pt}
\begin{tabular}{lcccccc}
\hline
\textbf{Scheme} &
\textbf{PQ} &
\textbf{Threshold} &
\textbf{Policy-Bound} &
\textbf{Cost/Overlap-} &
\textbf{State-Bound} &
\textbf{Reusable} \\
&
\textbf{Security} &
\textbf{Authorization} &
\textbf{Authorization} &
\textbf{Aware Redaction} &
\textbf{Execution} &
\textbf{Auditing} \\
\hline

Scheme~\cite{ref1}
& $\times$ & $\checkmark$ & $\times$ & $\times$
& $\times$ & $\times$ \\

Scheme~\cite{ref5}
& $\times$ & $\triangle$ & $\checkmark$ & $\times$
& $\triangle$ & $\times$ \\

Scheme~\cite{ref13}
& $\times$ & $\times$ & $\times$ & $\times$
& $\triangle$ & $\triangle$ \\

Scheme~\cite{ref14}
& $\times$ & $\times$ & $\triangle$ & $\triangle$
& $\triangle$ & $\times$ \\

Scheme~\cite{ref15}
& $\times$ & $\checkmark$ & $\checkmark$ & $\times$
& $\triangle$ & $\times$ \\

Scheme~\cite{ref17}
& $\checkmark$ & $\times$ & $\triangle$ & $\times$
& $\times$ & $\times$ \\

Scheme~\cite{ref20}
& $\times$ & $\triangle$ & $\checkmark$ & $\triangle$
& $\triangle$ & $\triangle$ \\

Scheme~\cite{ref27}
& $\times$ & $\checkmark$ & $\times$ & $\times$
& $\triangle$ & $\times$ \\

Scheme~\cite{ref33}
& $\times$ & $\triangle$ & $\checkmark$ & $\triangle$
& $\triangle$ & $\triangle$ \\

Scheme~\cite{ref34}
& $\times$ & $\checkmark$ & $\checkmark$ & $\times$
& $\triangle$ & $\times$ \\

\hline
\textbf{VeRedact-PQ}
& $\checkmark$ & $\checkmark$ & $\checkmark$
& $\checkmark$ & $\checkmark$ & $\checkmark$ \\
\hline

\multicolumn{7}{l}{
\footnotesize $\checkmark$: Supported;
$\triangle$: Partially supported;
$\times$: Not supported.}
\end{tabular}
\end{table*}
    \subsection{Distributed and Scalable Blockchain Redaction}
    
    Redactable blockchains commonly employ chameleon hashes to modify
    committed data while preserving blockchain consistency~\cite{ref9}.
    Subsequent studies extended this model to consortium blockchains
    ~\cite{ref2} and decentralized redaction through distributed chameleon
    hashes and trapdoor management~\cite{ref1,ref3,ref6,ref19,ref30}.
    Multi-party chameleon hashing~\cite{ref12} and robust threshold
    redaction~\cite{ref27} further reduce dependence on a single redaction
    authority, while dynamic and controlled constructions support evolving
    redaction requirements~\cite{ref5,ref16,ref36}.
    
    Scalability has also been addressed through efficient update
    mechanisms, lightweight storage, specialized Merkle structures, and
    application-oriented redaction~\cite{ref14,ref20,ref21,ref23,ref26}.
    However, efficiency is typically optimized at a particular layer,
    whereas high-volume redaction incurs combined costs from target lookup,
    request validation, committee authorization, authenticated state
    updates, and chameleon-hash adaptation. VeRedact-PQ addresses these
    costs jointly through sharded lookup, adaptive request batching,
    coalesced Merkle updates, and one distributed PQCH adaptation per
    affected transaction-batch root transition.
    
    \subsection{Policy-Controlled, Private, and Verifiable Redaction}
    
    Fine-grained redaction has been developed using attribute- and
    policy-based mechanisms~\cite{ref8,ref15}, while privacy-preserving and
    policy-hiding constructions restrict disclosure of sensitive
    authorization information~\cite{ref4,ref18,ref24,ref31,ref34}. Other schemes
    strengthen traceability and accountability through dynamic policies,
    anonymous but accountable redaction, policy-compliant rewriting, and
    cross-chain accountability~\cite{ref25,ref28,ref29,ref32,ref35}. Verifiable
    redaction has likewise been studied through efficient query and
    integrity auditing~\cite{ref13}, lightweight verification and
    permission supervision~\cite{ref20}, and publicly auditable
    architectures~\cite{ref33}.
    
    Despite these advances, policy validation, multi-party authorization,
    redaction execution, and auditing are often treated as separate
    functions. VeRedact-PQ instead binds them at request level: private
    policy eligibility is verified before batching, committee authorization
    is bound to an authenticated batch, and PBRP links this evidence to the
    executed state transition. A sharded RAI subsequently supports
    query-based auditing using compact multiproofs and shared batch
    evidence, preserving individual accountability without repeatedly
    processing complete redaction proofs.
    
    \subsection{Post-Quantum Redactable Blockchains}
    
    Post-quantum redactable blockchain research has primarily focused on
    the chameleon-hash primitive. Existing approaches include
    quantum-resistant key-exposure-free chameleon hashing~\cite{ref11},
    lattice-based redactable blockchains~\cite{ref10}, quantum-resistant
    redaction with trapdoor updates~\cite{ref17}, and lattice-based policy-
    or identity-oriented chameleon hashes~\cite{ref18,ref37}. These works
    provide important foundations for protecting controlled redaction
    against quantum adversaries.
    
    However, a quantum-resistant chameleon hash alone does not provide an
    end-to-end post-quantum redaction workflow if authentication,
    authorization, committee approval, or auditing still relies on
    classical mechanisms. VeRedact-PQ therefore extends PQ protection
    across the security-critical workflow using PQ authentication,
    PQZK-based private policy validation, PQ multi-party authorization, and
    distributed PQCH redaction. To control overhead, expensive PQZK
    verification is performed only when required, while compact PQ-signed
    attestations and hash-based authenticated proofs support subsequent
    authorization and routine auditing.
    
    Overall, VeRedact-PQ differs from existing work by jointly addressing
    \emph{post-quantum security, scalable multi-party redaction, and
    privacy-preserving verifiability} across the complete redaction
    lifecycle rather than optimizing these properties independently.
    

As summarized in Table~\ref{tab:comparison}, existing schemes typically
address only subsets of secure and scalable redaction. Decentralized
and threshold approaches strengthen distributed control
~\cite{ref1,ref15,ref27}, policy-oriented schemes improve authorization
~\cite{ref5,ref34}, and other works emphasize scalability
~\cite{ref14,ref20} or auditing~\cite{ref13,ref33}. Quantum-resistant
designs such as~\cite{ref17} protect the redaction primitive but do not
jointly provide post-quantum, policy-bound authorization, scalable
execution, and verification.

VeRedact-PQ integrates these properties across the redaction lifecycle.
RADC and state-bound PQZK bind authorization to authenticated state;
CORA exploits execution locality among validated requests; and SFMR
converts this locality into shared Merkle processing and batch-level
PQCH adaptation. PBRP binds authorization to the executed transition,
while SAVER enables state-bound reuse of shared audit evidence and
selective PQZK re-verification for challenged requests. Thus,
VeRedact-PQ improves scalability without relaxing request-level
authorization or state-transition accountability.


 \section{Our Proposed VeRedact-PQ Scheme}

This section presents the system model, threat model, and protocol of
VeRedact-PQ.

\subsection{System Model}

\begin{figure}[t]
\centering
\includegraphics[width=1\linewidth]{VeRedact-PQ System Model.png}
\caption{System model of the proposed VeRedact-PQ framework}
\label{fig:system_model}
\end{figure}

As illustrated in Fig.~\ref{fig:system_model}, VeRedact-PQ considers a
permissioned blockchain operated by mutually accountable consortium
organizations. Transactions may subsequently require legitimate
modification or removal under predefined redaction policies. Rather
than assigning redaction capability to a single authority, VeRedact-PQ
separates request validation, threshold authorization, redaction
execution, and post-redaction auditing. The system comprises the
following entities.

\textbf{1) Data Owner (DO):}
The DO generates and signs transactions before submission to the
blockchain, providing origin authentication and integrity. Each
transaction is associated with an identifier, data type, timestamp, and
redaction policy. A DO may also act as an authorized requester when
permitted by the applicable policy.

\textbf{2) Authorized Requester (AR):}
An AR is an authenticated entity permitted to request modification or
removal of committed data. It submits a PQ-signed request together with
the evidence required to satisfy the applicable public and private
policy conditions. An AR may be the original DO or another entity
granted redaction rights by the consortium.

\textbf{3) Redaction Committee (RC):}
The RC, denoted $\mathcal{C}_e$, consists of $n$ consortium members
appointed for authorization epoch $e$. At least $t$ members must approve
a redaction batch. Committee members also hold distributed shares of
the PQCH trapdoor and participate in authorized PQCH adaptation.
Consequently, fewer than $t$ committee members cannot independently
exercise the redaction capability.

\textbf{4) Validation and Policy Service (VPS):}
The VPS authenticates requesters, resolves authenticated transaction
state, evaluates public policy conditions, and verifies state-bound
PQZK proofs for private policy conditions. Successfully validated
requests receive PQ-signed validation attestations and enter the
pending queue $\mathcal{Q}_e$. The VPS also issues signed admission
receipts that provide verifiable request waiting times for subsequent
CORA processing.

\textbf{5) Permissioned Blockchain Network (PBN):}
The PBN comprises authenticated consortium nodes executing the
underlying permissioned consensus protocol. It maintains transaction
and policy commitments, Merkle-batch states, PQCH commitments,
authenticated index roots, authorization evidence, and finalized
redaction checkpoints. Large objects, including PQZK proofs and PBRP
evidence, may be retained in an off-ledger evidence store, with their
digests or authenticated roots anchored on-chain.

The PBN additionally provides two service functions. A
\emph{checkpoint authority} signs finalized batch/index checkpoints,
while an \emph{audit service} maintains the Redaction Audit Index (RAI),
resolves authorized audit queries, and signs responses. Their outputs
remain independently verifiable against blockchain-anchored
commitments.

\textbf{6) Auditor (AU):}
An AU is an authorized internal or external verifier that examines
completed redactions without requiring access to the underlying
sensitive transaction content. It verifies request validation,
committee authorization, redaction-state transitions, PQCH commitment
preservation, and finalized blockchain evidence at the verification
depth required by the audit.

\medskip
VeRedact-PQ organizes redaction around three complementary optimization
mechanisms. \emph{Cost- and Overlap-Aware Redaction Aggregation (CORA)}
selects already validated requests using authenticated batch/leaf
locality and bounded waiting time, reducing the expected number of
affected batches and Merkle-frontier nodes. \emph{Shared-Frontier
Merkle Redaction (SFMR)} exploits this locality during execution by
reusing common authentication paths and performing one PQCH adaptation
per affected transaction batch. The resulting authorization and
transition dependencies are preserved by the
\emph{Policy-Bound Batch Redaction Proof (PBRP)}.

Finally, \emph{State-Aware Verification Evidence Reuse (SAVER)}
organizes PBRP evidence into state-bound verification dependencies.
Shared authorization, transition, and checkpoint evidence can therefore
be verified once under an authenticated state context and reused across
related audit results, while challenged requests can selectively
escalate to independent PQZK verification.

Accordingly, VeRedact-PQ maintains request-level policy validation
before optimization, threshold committee authorization before
redaction, authorization-bound SFMR/PQCH execution, and state-aware
post-redaction verification. This separation prevents batching or
verification reuse from expanding redaction authority while allowing
cryptographic work to be shared across request selection, execution,
and auditing.
\subsection{Threat Model}

VeRedact-PQ considers malicious or compromised participants in a
permissioned blockchain and quantum-capable adversaries. The main
threats and assumptions are as follows.

\textbf{1) Malicious Requester:}
A requester may submit unauthorized, malformed, conflicting, or replayed
redaction requests, modify data outside its policy scope, or reuse
evidence from another state. Requests are therefore PQ-authenticated
and validated against the applicable policy and authenticated
transaction state before CORA processing.

\textbf{2) Malicious Committee Members:}
Committee members may collude to authorize invalid redactions or misuse
PQCH shares. We assume fewer than $t$ members of the active committee
$\mathcal{C}_e$ are compromised. PQCH shares are proactively refreshed
or reshared across epochs over authenticated confidential channels,
with obsolete shares securely erased. Security against cross-epoch
accumulation assumes fewer than $t$ contemporaneously valid shares are
compromised within each resharing interval.

\textbf{3) Malicious Blockchain or Service Participant:}
A compromised consortium node or audit service may alter, omit, inject,
or replay transaction, checkpoint, index, or audit evidence. We assume
the PBN satisfies its prescribed consensus fault-tolerance condition.
Authenticated roots, PQ signatures, PQCH commitments, and finalized
checkpoints make inconsistent evidence detectable. Audit completeness
is claimed only when authenticated coverage evidence is available.

\textbf{4) Authorization and Execution Manipulation:}
An adversary may insert or substitute requests after CORA selection,
reuse authorization for another state, or execute unauthorized
positions. The validated-request commitments, $R_e^{VR}$, $\Gamma_e$,
and $\Theta_b$ bind authorization to its execution context. Phase~5
enforces
\begin{equation}
\mathcal{P}_b^{exec}\subseteq\mathcal{P}_b^{auth},
\end{equation}
allowing stale requests to be withheld but preventing unauthorized
positions from being introduced.

\textbf{5) Verification-Reuse Adversary:}
An adversary may attempt to reuse SAVER evidence under a different
authorization, transition, or checkpoint state. SAVER binds each
verified result to its content and authenticated state context through
$\chi(v)$; any change to the evidence or its bound dependencies
requires fresh verification.

\textbf{6) Privacy-Curious Auditor:}
An auditor may attempt to infer information beyond its authorized
scope. PQZK hides private credentials and attributes, while
query-scoped evidence limits unnecessary content disclosure. Access-
pattern privacy and confidentiality against PBN nodes storing the
transaction content are outside the scope of this work.

\textbf{7) Quantum-Capable Adversary:}
The security-critical redaction path uses PQ signatures, PQZK, PQCH,
and quantum-resistant hash/PRF assumptions. The underlying permissioned
consensus and platform credentials are not required to be
post-quantum.

\medskip
\noindent\textbf{Trust and Scope Assumptions.}
The VPS is trusted to correctly authenticate requesters, evaluate
public policy conditions, and verify PQZK proofs before issuing
validation attestations. VPS compromise is outside the primary threat
model, although challenged audits can independently reverify the
underlying PQZK proof.

We assume the security of the underlying cryptographic primitives and
fewer than $t$ contemporaneously compromised committee shares.
Violation of the PBN consensus assumption, VPS or endpoint compromise,
denial-of-service, traffic-analysis leakage, and implementation-level
side-channel attacks are outside the scope of this work.
  \subsection{System Process}

VeRedact-PQ operates through six sequential phases covering system
initialization, transaction commitment, request validation, authorization,
redaction execution, and post-redaction auditing. The major notations
used throughout the framework are summarized in
Table~\ref{tab:notation}.

\begin{table*}[t]
\centering
\caption{Summary of Major Notations}
\label{tab:notation}
\scriptsize
\renewcommand{\arraystretch}{1.10}
\setlength{\tabcolsep}{3.8pt}
\begin{adjustbox}{max width=\textwidth}
\begin{tabular}{@{}ll@{\quad}ll@{}}
\hline
\textbf{Notation} & \textbf{Description} &
\textbf{Notation} & \textbf{Description}\\
\hline

$\lambda$, $PP$
& Security parameter, public parameters &
$R_i$, $\sigma_i^R$
& Redaction request, requester PQ signature\\

$e$, $e_P$, $e_i$, $e_b$, $e_A$
& Authorization, policy, state, batch, audit epochs &
$x_i$, $w_i$, $\pi_i^{PQ}$
& PQZK statement, witness, proof\\

$(pk_i,sk_i)$
& PQ signature key pair of entity $i$ &
$h_i^{\mathrm{st}}$
& Authenticated transaction-state digest\\

$sk_V$, $sk_{\mathrm{CP}}$, $sk_A$
& VPS, checkpoint, audit-service signing keys &
$C_i^{VR}$, $\alpha_i$
& Validated-request commitment, VPS attestation\\

$K_{\mathrm{idx}}$, $K_A$
& PRF keys for SA-RLI and RAI &
$VR_i$, $\mathcal{Q}_e$, $rc_i$
& Validated request, pending queue, admission receipt\\

$\mathcal{C}_e$, $n$, $t$
& Epoch committee, size, threshold &
$\mathcal{B}_e$, $\Omega_e$
& CORA-selected requests, affected batches\\

$pk_{\mathrm{CH}}$, $td_{e,k}$
& PQCH public key, trapdoor share &
$\eta_i$, $R_e^{VR}$
& Authorization leaf, validated-request root\\

$P_j$, $PID_j$, $v_j$, $C_{P_j}$
& Policy, identifier, version, commitment &
$\Gamma_e$
& CORA authorized execution-plan commitment\\

$m_i$, $\rho_i$
& Transaction content, commitment randomness &
$C_e^B$, $Auth_e^B$
& Authorization-batch commitment, evidence\\

$TID_i$, $DT_i$, $ts_i$
& Transaction identifier, data type, timestamp &
$\sigma_{e,k}^B$, $\mathcal{A}_e$
& Committee approval, approval set\\

$\sigma_i$, $C_i^{\mathrm{orig}}$
& Data-owner signature, provenance commitment &
$PBRP_i^{auth}$, $PBRP_i$
& Authorization and final redaction evidence\\

$I_i$, $D_i$, $L_i$
& RADC binding, content commitment, Merkle leaf &
$X_i^{\mathrm{red}}$, $D_i'$
& Redaction payload, proposed content commitment\\

$MR_b$, $r_b$, $CH_b$
& Batch root, PQCH randomness, commitment &
$\mathcal{P}_b^{auth}$, $\mathcal{P}_b^{exec}$
& Authorized and executed positions\\

$v_b$, $A_b$
& Batch version, finalized batch checkpoint &
$\mathcal{F}_b$
& SFMR external authentication frontier\\

$\tau_i$, $E_i$
& SA-RLI token, authenticated index entry &
$\Theta_b$
& Authorization-bound SFMR transition commitment\\

$R_{\mathrm{RLI}}$, $v_{\mathrm{RLI}}$
& SA-RLI root, snapshot version &
$\delta_{b,k}$
& Distributed PQCH adaptation share\\

$RR_i$
& Finalized redaction record &
$\Xi_b$
& Shared batch execution-evidence commitment\\

$\tau_i^A$, $E_i^A$
& RAI token, authenticated audit entry &
$R_{\mathrm{RAI}}$, $v_A$
& RAI root, audit snapshot version\\

$CP_{e_A}^A$
& Blockchain-anchored RAI checkpoint &
$Q_j^A$, $Resp_j^A$
& Audit query, authenticated response\\

$G_{Q_j}^{V}$
& SAVER verification-dependency DAG &
$id(v)$, $\kappa(v)$
& Evidence-node identifier, state context\\

$\chi(v)$, $VC[\chi(v)]$
& SAVER verification context, cached status &
$\mathcal{D}_i^{\mathrm{routine}},
 \mathcal{D}_i^{\mathrm{deep}}$
& Routine and deep verification dependencies\\

\hline
\end{tabular}
\end{adjustbox}
\end{table*}
    
    \subsubsection*{\textbf{Phase 1: System Initialization}}

This phase initializes the cryptographic parameters, registered
entities, redaction policies, and distributed redaction authority of
VeRedact-PQ. Let $\lambda$ denote the security parameter and $e$ the
current authorization epoch.

\textbf{Step 1: Global Parameter Setup.}
Given $1^\lambda$, the consortium initializes the cryptographic
primitives and publishes
\begin{equation}
\begin{aligned}
PP=(&\lambda,H,H_1,H_2,H_A,H_V,\mathsf{PQCH},\\
&\mathsf{PQSIG},\mathsf{PQZK},\mathsf{Merkle},\mathsf{PRF}),
\end{aligned}
\end{equation}
where $H$ is a post-quantum-secure hash function;
$H_1$, $H_2$, $H_A$, and $H_V$ are domain-separated instances of
$H$; $\mathsf{PQCH}$ is a post-quantum chameleon-hash scheme;
$\mathsf{PQSIG}$ is a post-quantum signature scheme;
$\mathsf{PQZK}$ is a post-quantum zero-knowledge proof system;
$\mathsf{Merkle}$ is an authenticated Merkle-tree construction; and
$\mathsf{PRF}$ is a quantum-secure pseudorandom function.
The additional domain $H_V$ is reserved for authenticated
verification-evidence commitments used by the state-aware verification
protocol in Phase~6.

\textbf{Step 2: Entity Registration and Key Generation.}
Each registered entity $U_i$, including data owners, authorized
requesters, committee members, the Validation and Policy Service
(VPS), PBN checkpoint and audit services, and auditors, generates
\begin{equation}
(pk_i,sk_i)\leftarrow
\mathsf{PQSIG.KeyGen}(1^\lambda).
\end{equation}
The tuple $(ID_i,Role_i,pk_i)$ is recorded in the consortium membership
registry, while $sk_i$ remains private. In particular, the VPS uses
$(pk_V,sk_V)$ for validation attestations; the checkpoint authority uses
$(pk_{\mathrm{CP}},sk_{\mathrm{CP}})$ for checkpoint signatures; and the
audit service uses $(pk_A,sk_A)$ for authenticated audit responses.
A committee member $C_k$ active in epoch $e$ uses its registered pair
$(pk_{e,k},sk_{e,k})$.

The consortium additionally samples
\begin{equation}
K_{\mathrm{idx}},K_A
\stackrel{\$}{\leftarrow}\{0,1\}^{\lambda},
\end{equation}
where $K_{\mathrm{idx}}$ derives SA-RLI lookup tokens and is available
to the PBN and VPS, whereas $K_A$ derives RAI audit tokens and is held
by the PBN audit service. The parameter $\lambda$ is selected according
to the target post-quantum security level.

\textbf{Step 3: Redaction Policy Registration.}
For each supported transaction or data class, the consortium defines a
redaction policy $P_j$ specifying eligible requester roles, permitted
operations, applicable conditions, and validity period. Its authenticated
state is committed as
\begin{equation}
C_{P_j}
=
H(PID_j\parallel P_j\parallel v_j\parallel e_P),
\end{equation}
where $PID_j$, $v_j$, and $e_P$ denote the policy identifier, version,
and commitment epoch, respectively. The tuple
$(PID_j,C_{P_j},v_j,e_P)$ is registered on-chain and subsequently binds
request validation and redaction authorization to an authenticated
policy state.

\textbf{Step 4: Epoch Committee Formation.}
At the beginning of epoch $e$, the consortium establishes an
$n$-member redaction committee
\begin{equation}
\mathcal{C}_e=(C_1,C_2,\ldots,C_n),\qquad t\leq n,
\end{equation}
where at least $t$ members are required to authorize and exercise the
distributed redaction capability. The configuration
$(e,\mathcal{C}_e,t)$ is recorded on-chain, binding subsequent
authorization and PQCH adaptation to the committee active in epoch $e$.

\textbf{Step 5: Distributed PQCH Trapdoor Setup.}
At system initialization, the initial committee
$\mathcal{C}_{e_0}$ jointly executes
\begin{equation}
(pk_{\mathrm{CH}},\{td_{e_0,k}\}_{k=1}^{n})
\leftarrow
\mathsf{PQCH.DKeyGen}(PP,\mathcal{C}_{e_0},t),
\end{equation}
where $pk_{\mathrm{CH}}$ is the public chameleon-hash key and
$td_{e_0,k}$ is the trapdoor share of member $C_k$.
The public key $pk_{\mathrm{CH}}$ is recorded on the PBN and remains
fixed across authorization epochs.

At each transition $e-1\rightarrow e$, the trapdoor is proactively
reshared to the incoming committee:
\begin{equation}
\{td_{e,k}\}_{k=1}^{n}
\leftarrow
\mathsf{PQCH.Reshare}
(\{td_{e-1,k}\},\mathcal{C}_e,t).
\end{equation}
The resharing protocol is executed over authenticated confidential
channels, and previous-epoch shares are securely erased after successful
resharing. Hence, the active committee can adapt existing batch
commitments without reconstructing the complete trapdoor at any single
member. Security against cross-epoch share accumulation assumes secure
erasure and that fewer than $t$ valid shares are compromised within
each resharing interval.

\textbf{Step 6: Post-Quantum ZK Parameter Setup.}
For the policy-compliance relation $\mathcal{R}_P$, the consortium
initializes
\begin{equation}
(pp_{\mathrm{PQZK}},pk_{\mathrm{PQZK}},vk_{\mathrm{PQZK}})
\leftarrow
\mathsf{PQZK.Setup}(1^\lambda,\mathcal{R}_P).
\end{equation}
The selected PQZK construction provides completeness, soundness, and
zero knowledge against quantum-capable adversaries. A transparent,
trapdoor-free setup is assumed; therefore, the public proving and
verification parameters can be published on the PBN without requiring
trusted setup-secret erasure or regeneration when new policy instances
are registered. Any prover-specific secret state required by the
selected instantiation remains with the corresponding prover.

All public keys, policy commitments, committee configurations, and
other security-critical initialization records are submitted as
transactions authenticated by the registering consortium members'
post-quantum signatures. Consequently, the cryptographic redaction
control plane does not rely solely on the underlying PBN's platform
credentials.
    
    
    
   \subsubsection*{\textbf{Phase 2: Redaction-Aware Transaction Commitment}}

This phase transforms each accepted transaction into a redaction-aware
authenticated state. VeRedact-PQ separates immutable transaction and
policy context from mutable content, commits transactions in
PQCH-protected Merkle batches, and constructs an authenticated lookup
index that exposes the batch and leaf locality required by subsequent
authorization and optimized redaction. Unless stated otherwise, the
operations in this phase are performed by the PBN.

\medskip
\noindent\textbf{Step 1: Transaction Authentication and Immutable
Provenance.}
A registered data owner $DO_i$ submits content $m_i$ with transaction
identifier $TID_i$, data type $DT_i$, redaction policy $PID_i$, and
timestamp $ts_i$. The owner samples
$\rho_i\stackrel{\$}{\leftarrow}\{0,1\}^{\lambda}$ and computes
\begin{equation}
D_i=H(m_i\parallel\rho_i),
\end{equation}
followed by
\begin{equation}
\begin{aligned}
\sigma_i\leftarrow
\mathsf{PQSIG.Sign}\big(
sk_{DO_i},
H(&TID_i\parallel D_i\parallel DT_i\\
&\parallel PID_i\parallel ts_i)\big).
\end{aligned}
\end{equation}
The submitted transaction is
\begin{equation}
TX_i=(TID_i,m_i,\rho_i,DT_i,PID_i,ts_i,\sigma_i).
\end{equation}
The PBN recomputes $D_i$ and verifies transaction uniqueness, the
registered owner key, $\sigma_i$, and the referenced policy before
acceptance. The verified key determines $ID_{DO_i}$.

To preserve the authenticated origin independently of later content
states, the PBN derives
\begin{equation}
C_i^{\mathrm{orig}}
=
H(TID_i\parallel D_i\parallel\sigma_i)
\end{equation}
and anchors the non-redactable provenance tuple
\begin{equation}
\mathcal{P}_i^{\mathrm{orig}}
=
(TID_i,DT_i,PID_i,ts_i,D_i,\sigma_i,C_i^{\mathrm{orig}}).
\end{equation}
Hence, $\sigma_i$ continues to authenticate the originally submitted
content commitment after an authorized redaction; subsequent content
states are authenticated by the redaction evidence generated in
Phases~4--6 rather than by re-signing them as original owner content.

\medskip
\noindent\textbf{Step 2: Redaction-Aware Dual Commitment.}
For each accepted transaction, the PBN samples a fresh nonce $n_i$ and
computes
\begin{equation}
Tag_i=H(TID_i\parallel n_i\parallel ts_i).
\end{equation}
The Redaction-Aware Dual Commitment (RADC) separates immutable context
from mutable content through
\begin{equation}
\begin{aligned}
I_i=H(&TID_i\parallel ID_{DO_i}\parallel DT_i
\parallel PID_i\\
&\parallel ts_i\parallel Tag_i)
\end{aligned}
\end{equation}
and the content commitment $D_i$ from Step~1. The corresponding
authenticated leaf is
\begin{equation}
L_i=H(I_i\parallel D_i).
\end{equation}

An authorized transition from $(m_i,\rho_i)$ to
$(m_i',\rho_i')$ therefore changes only
\begin{equation}
D_i\rightarrow D_i'=H(m_i'\parallel\rho_i'),
\end{equation}
while $I_i$ remains invariant. This separation prevents a content
redaction from altering the transaction identity, ownership context,
data type, policy binding, or original provenance. The current content
and opening randomness are maintained in the transaction state, whereas
the immutable provenance record remains unchanged.

\medskip
\noindent\textbf{Step 3: Redaction-Locality Batch Commitment.}
Authenticated leaves are placed at deterministic positions in a
transaction batch
\begin{equation}
\mathcal{T}_b=(L_{b,1},\ldots,L_{b,N_b}),
\end{equation}
whose root is
\begin{equation}
MR_b=
\mathsf{Merkle.Root}(L_{b,1},\ldots,L_{b,N_b}).
\end{equation}
The batch is then committed using
\begin{equation}
CH_b=
\mathsf{PQCH.Hash}(pk_{\mathrm{CH}},MR_b,r_b),
\qquad
r_b\stackrel{\$}{\leftarrow}\mathcal{R}_{\mathrm{CH}},
\end{equation}
where $\mathcal{R}_{\mathrm{CH}}$ denotes the randomness space of the
selected PQCH scheme. The initial batch version is $v_b=0$.

The Merkle topology is retained as authenticated update state. For any
future set of leaves $\mathcal{L}_b^{*}$ selected for redaction, define
its minimal external authentication frontier as
\begin{equation}
\mathcal{F}_b(\mathcal{L}_b^{*})
=
\left\{
u:
u\notin\mathsf{Anc}(\mathcal{L}_b^{*}),
\;
\mathsf{parent}(u)\in
\mathsf{Anc}(\mathcal{L}_b^{*})
\right\},
\end{equation}
where $\mathsf{Anc}(\mathcal{L}_b^{*})$ is the union of Merkle nodes on
paths from the selected leaves to $MR_b$. The frontier is not computed
during initial commitment; rather, the stored topology permits it to be
derived when multiple authorized redactions are coalesced in Phase~5.
Consequently, overlapping authentication paths can later be represented
once instead of independently for every modified leaf.

\begin{figure}[t]
\centering
\includegraphics[width=0.7\linewidth]{Merkle.PNG}
\caption{Redaction-aware Merkle batch and PQCH commitment structure.}
\label{fig:merkle}
\end{figure}

\medskip
\noindent\textbf{Step 4: Authenticated Redaction-State Indexing.}
To resolve redaction targets without scanning blockchain transactions,
VeRedact-PQ constructs the Sharded Authenticated Redaction Lookup Index
(SA-RLI). For each transaction,
\begin{equation}
\tau_i=\mathsf{PRF}_{K_{\mathrm{idx}}}(TID_i),
\end{equation}
and the corresponding authenticated state entry is
\begin{equation}
E_i=
(\tau_i,Tag_i,b,pos_i,PID_i,e_i,I_i,D_i,ptr_i),
\end{equation}
where $b$ and $pos_i$ identify the transaction's batch and leaf
position, $e_i$ identifies the epoch of the indexed state, and $ptr_i$
references the authenticated transaction state required to resolve the
leaf. The batch version $v_b$ is obtained from the finalized checkpoint
rather than duplicated in every entry.

Each entry is deterministically assigned to a shard and logical bucket:
\begin{equation}
sid_i=H_1(\tau_i)\bmod S,
\qquad
bid_i=H_2(\tau_i)\bmod B,
\end{equation}
where $S$ and $B$ are the numbers of index shards and logical buckets.
Within each shard, entries are canonically ordered by
$(bid_i,\tau_i)$ and committed as
\begin{equation}
R_s=
\mathsf{Merkle.Root}(E_{s,1},\ldots,E_{s,n_s}),
\end{equation}
with global authenticated index state
\begin{equation}
R_{\mathrm{RLI}}
=
\mathsf{Merkle.Root}(R_1,\ldots,R_S).
\end{equation}

The explicit $(b,pos_i)$ mapping is also used by the later
cost- and overlap-aware redaction protocol to identify requests that
share a transaction batch or overlapping Merkle paths without exposing
or modifying their authenticated content state.

\medskip
\noindent\textbf{Step 5: Atomic State Checkpointing.}
After the transaction batch and its SA-RLI updates are finalized, the
PBN assigns a monotonically increasing snapshot identifier
$v_{\mathrm{RLI}}$ and constructs
\begin{equation}
\begin{aligned}
A_b=\big(&b,CH_b,MR_b,r_b,R_{\mathrm{RLI}},
v_{\mathrm{RLI}},\\
& e,v_b,ts_b,\sigma_b\big),
\end{aligned}
\end{equation}
where $e$ is the checkpoint-finalization epoch. The checkpoint authority
computes
\begin{equation}
\begin{aligned}
\sigma_b=
\mathsf{PQSIG.Sign}\big(
sk_{\mathrm{CP}},
H(&b\parallel CH_b\parallel MR_b\parallel r_b\\
&\parallel R_{\mathrm{RLI}}\parallel v_{\mathrm{RLI}}
\parallel e\\
&\parallel v_b\parallel ts_b)\big).
\end{aligned}
\end{equation}
The signed checkpoint is anchored to the PBN, atomically binding the
PQCH-protected transaction-batch state to the corresponding finalized
SA-RLI snapshot. Thus, later authorization and redaction operate against
a uniquely identified tuple
\begin{equation}
\mathsf{State}_b=
(b,MR_b,CH_b,v_b,R_{\mathrm{RLI}},v_{\mathrm{RLI}},e),
\end{equation}
which prevents evidence from different batch or index versions from
being silently combined.

For lookup acceleration, each shard may maintain a Cuckoo filter
$CF_s$ over its current tokens. The filter is used only as a screening
structure: positive results require authenticated SA-RLI verification,
and a negative result is accepted only when $CF_s$ is synchronized with
the selected finalized snapshot; otherwise, authenticated lookup is
performed.

At completion of this phase, SA-RLI resolves each target to its
authenticated RADC state, policy, batch, and leaf position. The
$(b,pos_i)$ locality information subsequently enables cost-aware request
coalescing, while the retained Merkle topology enables shared-frontier
redaction without weakening the authenticated state binding.
    
    
   \subsubsection*{\textbf{Phase 3: Redaction Request Authentication and
PQZK-Based Policy Verification}}

This phase converts a redaction request into an authenticated,
policy-compliant, and state-bound validated request. VeRedact-PQ first
resolves the target against a finalized SA-RLI snapshot, authenticates
the requester, rejects requests that fail public policy checks, and
invokes PQZK only for the remaining private policy conditions. Accepted
requests receive a PQ-signed validation attestation and enter the pending
queue for cost- and overlap-aware authorization in Phase~4.

\medskip
\noindent\textbf{Step 1: Request Admission and Authenticated State
Resolution.}
A requester $U_r$ prepares
\begin{equation}
R_i=(ID_r,TID_i,op_i,D_i',e,ts_r,n_r),
\end{equation}
where $op_i$ is the requested redaction operation, $e$ is the
authorization epoch, and $(ts_r,n_r)$ provides freshness. The
corresponding redaction payload is
\begin{equation}
X_i^{\mathrm{red}}=(m_i',\rho_i'),
\qquad
D_i'=H(m_i'\parallel\rho_i'),
\end{equation}
so that $R_i$ commits to the proposed content state without exposing
$m_i'$ in the authorization record. The payload is retained for
execution in Phase~5, where its opening against $D_i'$ is verified
before any state transition.

After basic format, timestamp, nonce, and epoch screening, the VPS
derives
\begin{equation}
\tau_i=\mathsf{PRF}_{K_{\mathrm{idx}}}(TID_i),
\qquad
sid_i=H_1(\tau_i)\bmod S,
\qquad
bid_i=H_2(\tau_i)\bmod B.
\end{equation}
A synchronized Cuckoo filter may reject an absent token; otherwise, the
VPS performs authenticated SA-RLI lookup and retrieves
\begin{equation}
E_i=(\tau_i,Tag_i,b,pos_i,PID_i,e_i,I_i,D_i,ptr_i).
\end{equation}

The VPS verifies $E_i$ against the corresponding shard root and
$R_{\mathrm{RLI}}$, retrieves the signed checkpoint $A_b$, and verifies
\begin{equation}
L_i=H(I_i\parallel D_i)
\end{equation}
at position $pos_i$ against $MR_b$. It also verifies the checkpoint
signature and consistency between the batch state and finalized index
snapshot. The resolved state is represented by
\begin{equation}
\mathsf{State}_b=
(b,MR_b,CH_b,v_b,R_{\mathrm{RLI}},v_{\mathrm{RLI}},e_b),
\end{equation}
where $e_b$ denotes the checkpoint-finalization epoch, distinguished
from the current authorization epoch $e$. The indexed epoch $e_i$
remains transaction-state metadata and need not equal either $e_b$ or
$e$.

The pair $(b,pos_i)$ is retained as authenticated locality metadata.
It is used only after validation to estimate common-batch and
Merkle-path overlap in Phase~4 and does not affect request eligibility
or policy satisfaction.

\medskip
\noindent\textbf{Step 2: Post-Quantum Requester Authentication.}
The requester signs the complete request:
\begin{equation}
\sigma_i^R\leftarrow
\mathsf{PQSIG.Sign}(sk_r,H(R_i))
\end{equation}
and submits $(R_i,\sigma_i^R)$. Using the registered key $pk_r$, the
VPS verifies
\begin{equation}
\mathsf{PQSIG.Verify}
(pk_r,H(R_i),\sigma_i^R)=1,
\end{equation}
together with the binding between $ID_r$ and $pk_r$, requester
registration, and eligibility to submit requests in epoch $e$.
The nonce $n_r$ is consumed in the authenticated requester's replay
domain only after successful authentication.

Only after authentication does the VPS release to the requester the
authenticated policy and state information required to construct the
PQZK statement.

\medskip
\noindent\textbf{Step 3: Public Policy and State Validation.}
Using the authenticated $PID_i$, the VPS retrieves policy $P_i$, its
version $v_i$, commitment epoch $e_P$, and on-chain commitment
$C_{P_i}$, and verifies
\begin{equation}
C_{P_i}
\stackrel{?}{=}
H(PID_i\parallel P_i\parallel v_i\parallel e_P).
\end{equation}
The VPS applies the registered policy-version and epoch-transition rules
to confirm that $P_i$ governs the resolved transaction state and is
valid for authorization epoch $e$. It then verifies $op_i$, all public
policy predicates, request freshness, and the validity of the resolved
state $\mathsf{State}_b$.

Requests failing these checks are rejected before PQZK verification.
This staged validation avoids invoking the comparatively expensive
PQZK verifier for requests that already fail authenticated public
conditions.

\medskip
\noindent\textbf{Step 4: State-Bound PQZK Policy Verification.}
For requests that pass the public checks, the requester proves
satisfaction of the private policy predicates without revealing its
credentials or sensitive authorization attributes. To bind the proof
to the exact authenticated state, define
\begin{equation}
h_i^{\mathrm{st}}
=
H\big(
TID_i\parallel I_i\parallel D_i\parallel pos_i
\parallel \mathsf{enc}(\mathsf{State}_b)
\big),
\end{equation}
where $\mathsf{enc}(\cdot)$ denotes an unambiguous canonical encoding.
The public statement is
\begin{equation}
\begin{aligned}
x_i=\big(&ID_r,TID_i,op_i,D_i',
C_{P_i},h_i^{\mathrm{st}},\\
& e_i,e,ts_r,n_r\big),
\end{aligned}
\end{equation}
and the witness $w_i$ contains the private credentials,
authorization attributes, and policy evidence.

The policy relation is
\begin{equation}
\begin{aligned}
\mathcal{R}_{P}(x_i,w_i)=1
\iff{}&
\mathsf{ValidCred}(w_i,C_{P_i},e)\\
&\land
\mathsf{RequesterBound}(w_i,ID_r)\\
&\land
\mathsf{PrivatePolicy}(x_i,w_i),
\end{aligned}
\end{equation}
where credential validity includes the applicable expiry and revocation
conditions, and requester binding proves credential ownership or valid
delegation.

The requester computes
\begin{equation}
\pi_i^{PQ}\leftarrow
\mathsf{PQZK.Prove}
(pk_{\mathrm{PQZK}},x_i,w_i),
\end{equation}
and the VPS accepts only if
\begin{equation}
\mathsf{PQZK.Verify}
(vk_{\mathrm{PQZK}},x_i,\pi_i^{PQ})=1.
\end{equation}
The state digest $h_i^{\mathrm{st}}$ prevents a valid proof from being
transplanted to a different transaction position, content state, batch
version, or finalized SA-RLI snapshot.

\medskip
\noindent\textbf{Step 5: Validated-Request Attestation and Queue
Admission.}
For each accepted request, the VPS assigns a unique $RID_i$ and derives
the proof digest
\begin{equation}
h_i^{PQ}=H(\pi_i^{PQ}).
\end{equation}
It then constructs the validated-request commitment
\begin{equation}
\begin{aligned}
C_i^{VR}
=
H\big(&RID_i\parallel H(R_i)
\parallel h_i^{\mathrm{st}}
\parallel C_{P_i}\\
&\parallel h_i^{PQ}
\parallel e\big)
\end{aligned}
\end{equation}
and signs
\begin{equation}
\alpha_i=
\mathsf{PQSIG.Sign}
\left(
sk_V,
H(RID_i\parallel C_i^{VR}\parallel e)
\right).
\end{equation}

The compact validated request is
\begin{equation}
\begin{aligned}
VR_i=\big(&RID_i,C_i^{VR},b,pos_i,PID_i,
v_b,v_{\mathrm{RLI}},\\
&e_i,e,h_i^{\mathrm{st}},h_i^{PQ},
\sigma_i^R,\alpha_i\big).
\end{aligned}
\end{equation}
The explicit $(b,pos_i)$ fields provide the authenticated locality
metadata required by Phase~4, whereas
$(v_b,v_{\mathrm{RLI}},h_i^{\mathrm{st}})$ bind that metadata to the
validated state. The complete request $R_i$, statement $x_i$, proof
$\pi_i^{PQ}$, policy evidence, and authenticated transaction evidence
remain available for independent or challenged verification.

The validated request is admitted to the pending queue
$\mathcal{Q}_e$ at time $ts_i^{\mathrm{adm}}$. The VPS returns
\begin{equation}
rc_i=
\mathsf{PQSIG.Sign}
\left(
sk_V,
H(RID_i\parallel C_i^{VR}\parallel
ts_i^{\mathrm{adm}})
\right),
\end{equation}
which provides verifiable evidence of admission. A request that reaches
the maximum waiting bound $T_{\max}$ becomes mandatory for the next
eligible authorization round, subject to state and policy revalidation.

Phase~4 operates only on such validated requests. It verifies
$\alpha_i$ and current-state freshness before applying
cost- and overlap-aware redaction aggregation; locality therefore
affects batching efficiency but never substitutes for requester,
policy, or state authorization.
    
    
  \subsubsection*{\textbf{Phase 4: Cost- and Overlap-Aware Redaction
Aggregation and Multi-Party Authorization}}

This phase applies \emph{Cost- and Overlap-Aware Redaction Aggregation
(CORA)} to validated requests. Unlike load-only batching, CORA exploits
authenticated batch and leaf locality to co-schedule requests that can
share Merkle-update work and PQCH adaptation, while enforcing a bounded
waiting time. The active committee then revalidates the selected state
and issues $t$-of-$n$ post-quantum authorization over the resulting
batch commitment.

\medskip
\noindent\textbf{Step 1: CORA Candidate Selection.}
Let $\mathcal{Q}_e$ denote the queue of Phase~3 validated requests for
authorization epoch $e$. For each $VR_i\in\mathcal{Q}_e$, the authenticated
pair $(b_i,pos_i)$ identifies its target transaction batch and Merkle
position. For a candidate set $\mathcal{S}\subseteq\mathcal{Q}_e$, define
\begin{equation}
\Omega(\mathcal{S})
=
\{b_i:VR_i\in\mathcal{S}\},
\end{equation}
and, for each $b\in\Omega(\mathcal{S})$,
\begin{equation}
\mathcal{L}_b(\mathcal{S})
=
\{pos_i:VR_i\in\mathcal{S},\,b_i=b\}.
\end{equation}
Using the fixed Merkle topology from Phase~2, the VPS derives the
corresponding external frontier
$\mathcal{F}_b(\mathcal{L}_b(\mathcal{S}))$. CORA estimates the
redaction cost as
\begin{equation}
\mathcal{C}_{\mathrm{CORA}}(\mathcal{S})
=
w_{\mathrm{CH}}|\Omega(\mathcal{S})|
+
w_{\mathrm{F}}
\sum_{b\in\Omega(\mathcal{S})}
|\mathcal{F}_b(\mathcal{L}_b(\mathcal{S}))|,
\end{equation}
where $w_{\mathrm{CH}}$ and $w_{\mathrm{F}}$ reflect the relative cost
of a distributed PQCH adaptation and an authenticated frontier node,
respectively.

Given a workload-dependent target size
\begin{equation}
\widehat{B}_e
=
\min\!\left\{
B_{\max},
\max\!\left\{
B_{\min},
\left\lceil
f(\lambda_e,|\mathcal{Q}_e|)
\right\rceil
\right\}
\right\},
\end{equation}
CORA selects
\begin{equation}
\mathcal{B}_e^{\mathrm{cand}}
=
\arg\min_{\substack{
\mathcal{S}\subseteq\mathcal{Q}_e\\
|\mathcal{S}|\leq\widehat{B}_e\\
\mathcal{M}_e\subseteq\mathcal{S}
}}
\mathcal{C}_{\mathrm{CORA}}(\mathcal{S}),
\end{equation}
where
\begin{equation}
\mathcal{M}_e=
\{VR_i\in\mathcal{Q}_e:
t_{\mathrm{now}}-ts_i^{\mathrm{adm}}\geq T_{\max}\}
\end{equation}
is the mandatory set of requests reaching the waiting bound. If
$|\mathcal{M}_e|>\widehat{B}_e$, the oldest mandatory requests are
processed first and the target expands up to $B_{\max}$ when capacity
permits. Thus, locality can improve execution efficiency but cannot
indefinitely postpone an admitted request.

The combinatorial objective need not be solved exhaustively. The
implementation uses a deterministic greedy realization that admits
mandatory requests first and then repeatedly selects the request with
the smallest marginal cost
\begin{equation}
\Delta_i(\mathcal{S})
=
\mathcal{C}_{\mathrm{CORA}}
(\mathcal{S}\cup\{VR_i\})
-
\mathcal{C}_{\mathrm{CORA}}(\mathcal{S}),
\end{equation}
with admission time and $RID_i$ used as deterministic tie-breakers.

\medskip
\noindent\textbf{Step 2: Attestation and Current-State Revalidation.}
Before authorization, each selected request is independently
revalidated. The committee verifies
\begin{equation}
\mathsf{PQSIG.Verify}
\left(
pk_V,
H(RID_i\parallel C_i^{VR}\parallel e),
\alpha_i
\right)=1
\end{equation}
and reconstructs $C_i^{VR}$ from the retained Phase~3 evidence.
Routine authorization verifies the VPS attestation and proof digest
rather than repeating the underlying PQZK verification.

The committee then resolves the latest authenticated state and
recomputes
\begin{equation}
\widehat{h}_i^{\mathrm{st}}
=
H\big(
TID_i\parallel I_i\parallel D_i^{cur}\parallel pos_i
\parallel
\mathsf{enc}(\mathsf{State}_b^{cur})
\big).
\end{equation}
The request remains eligible only if
\begin{equation}
\widehat{h}_i^{\mathrm{st}}
=
h_i^{\mathrm{st}},
\end{equation}
epoch $e$ remains active, and the attested policy version remains
applicable. Requests that fail revalidation are removed and returned
for Phase~3 revalidation. CORA is not rerun merely to replace removed
requests; this prevents post-selection state changes from altering the
meaning of the candidate authorization round.

Let the remaining eligible requests, canonically ordered by
$(b_i,pos_i,RID_i)$, be
\begin{equation}
\mathcal{B}_e=(VR_1,\ldots,VR_m).
\end{equation}

\medskip
\noindent\textbf{Step 3: State-Bound Authorization Batch Commitment.}
For each $VR_i\in\mathcal{B}_e$, define
\begin{equation}
\eta_i
=
H_V\big(
\texttt{AUTH-LEAF}\parallel
RID_i\parallel C_i^{VR}\parallel
\alpha_i
\big),
\end{equation}
where the domain-separated $H_V$ from Phase~1 prevents
cross-context reuse. The leaves are committed as
\begin{equation}
R_e^{VR}
=
\mathsf{Merkle.Root}
(\eta_1,\ldots,\eta_m).
\end{equation}
Because $C_i^{VR}$ already binds the request, policy commitment,
PQZK-proof digest, authorization epoch, and authenticated state digest,
$\eta_i$ need not duplicate these fields.

To bind authorization to the execution structure selected by CORA,
define
\begin{equation}
\Gamma_e=
H_V\!\left(
\texttt{CORA-PLAN}\parallel
\mathsf{enc}
\left(
\{(b_i,pos_i)\}_{i=1}^{m}
\right)
\right),
\end{equation}
and construct
\begin{equation}
C_e^{B}
=
H_V\big(
\texttt{AUTH-BATCH}\parallel
BID_e\parallel e\parallel m\parallel
R_e^{VR}\parallel\Gamma_e\parallel ts_e
\big),
\end{equation}
where $BID_e$ uniquely identifies the authorization batch.
The plan commitment $\Gamma_e$ fixes the authorized batch/leaf
locality without committing to a particular execution implementation.

\medskip
\noindent\textbf{Step 4: Multi-Party Committee Authorization.}
Each approving member $C_k\in\mathcal{C}_e$ independently verifies the
finalized request set, current-state checks, $R_e^{VR}$, and
$\Gamma_e$, and signs
\begin{equation}
\sigma_{e,k}^{B}
\leftarrow
\mathsf{PQSIG.Sign}
\left(
sk_{e,k},C_e^{B}
\right).
\end{equation}
The valid approval set is
\begin{equation}
\begin{aligned}
\mathcal{A}_e=\big\{&
(k,\sigma_{e,k}^{B}):
C_k\in\mathcal{C}_e,\\
&
\mathsf{PQSIG.Verify}
(pk_{e,k},C_e^{B},\sigma_{e,k}^{B})=1
\big\}.
\end{aligned}
\end{equation}
Authorization succeeds only when approvals from at least $t$ distinct
registered committee members are present:
\begin{equation}
|\mathcal{A}_e|\geq t.
\end{equation}
The resulting shared authorization evidence is
\begin{equation}
Auth_e^{B}
=
(BID_e,C_e^{B},R_e^{VR},\Gamma_e,
\mathcal{A}_e,e,ts_e).
\end{equation}
This construction uses independent post-quantum signatures and does not
require a native threshold-signature scheme.

\medskip
\noindent\textbf{Step 5: Policy-Bound Authorization Evidence.}
For each authorized request, the system derives a Merkle membership
proof $MP_i^{B}$ for $\eta_i$ under $R_e^{VR}$ and constructs
\begin{equation}
PBRP_i^{auth}
=
(RID_i,\eta_i,MP_i^{B},Auth_e^{B}).
\end{equation}
Since $Auth_e^{B}$ already contains $C_e^{B}$ and $R_e^{VR}$, they are
not duplicated in each authorization component. Shared committee
evidence can therefore be verified once per authorization batch, while
$MP_i^{B}$ binds each request to that evidence.

Only threshold-authorized batches proceed to Phase~5. Before producing
a PQCH adaptation share, each committee member verifies that the
proposed state transition contains only requests committed by
$Auth_e^{B}$ and that its affected $(b_i,pos_i)$ set agrees with
$\Gamma_e$. Phase~5 then applies shared-frontier Merkle redaction to the
authorized positions and binds the resulting state transition and PQCH
adaptation to $PBRP_i^{auth}$.
    
 \subsubsection*{\textbf{Phase 5: Shared-Frontier Redaction and
Verifiable State Transition}}

This phase executes the CORA-authorized requests and converts their
shared locality into cryptographic execution savings. For each affected
transaction batch, VeRedact-PQ applies \emph{Shared-Frontier Merkle
Redaction (SFMR)} to recompute common Merkle paths once, performs one
threshold PQCH adaptation for the resulting batch-root transition, and
atomically binds the new transaction and SA-RLI states to the Phase~4
authorization.

\medskip
\noindent\textbf{Step 1: Authorization-Bound Execution Validation.}
Given the authorized request set $\mathcal{B}_e$, the PBN verifies
$Auth_e^B$, the threshold approval set $\mathcal{A}_e$, and each
$PBRP_i^{auth}$. It reconstructs the authorized position set and verifies
\begin{equation}
\Gamma_e
\stackrel{?}{=}
H_V\!\left(
\texttt{CORA-PLAN}\parallel
\mathsf{enc}\big(\{(b_i,pos_i)\}_{VR_i\in\mathcal{B}_e}\big)
\right).
\end{equation}
The requests are then partitioned by target batch:
\begin{equation}
\mathcal{B}_e
=
\bigcup_{b\in\Omega_e}\mathcal{B}_{e,b},
\qquad
\mathcal{B}_{e,b}
=
\{VR_i\in\mathcal{B}_e:b_i=b\},
\end{equation}
where $\Omega_e$ is the set of affected transaction batches.

Immediately before execution, the PBN resolves the latest authenticated
state for every request and recomputes
\begin{equation}
\widehat{h}_i^{\mathrm{st}}
=
H\big(
TID_i\parallel I_i\parallel D_i^{cur}\parallel pos_i
\parallel
\mathsf{enc}(\mathsf{State}_b^{cur})
\big).
\end{equation}
Execution requires
\begin{equation}
\widehat{h}_i^{\mathrm{st}}=h_i^{\mathrm{st}},
\end{equation}
together with a still-valid authorization epoch, policy, and approved
operation. Requests targeting the same transaction are conflict-checked;
at most one resulting content state is executed for a transaction in
one authorization round. Any stale, conflicting, or otherwise invalid
request causes the corresponding transition to be withheld rather than
silently replaced by an unauthorized request.

\medskip
\noindent\textbf{Step 2: RADC Transition Construction.}
For each executable $VR_i$, the PBN retrieves
\begin{equation}
X_i^{\mathrm{red}}=(m_i',\rho_i')
\end{equation}
and verifies
\begin{equation}
H(m_i'\parallel\rho_i')
\stackrel{?}{=}D_i'.
\end{equation}
A valid request induces the RADC transition
\begin{equation}
(I_i,D_i)\longrightarrow(I_i,D_i'),
\end{equation}
with
\begin{equation}
L_i=H(I_i\parallel D_i),
\qquad
L_i'=H(I_i\parallel D_i').
\end{equation}
Thus, execution modifies only the authorized content commitment while
preserving the immutable transaction-policy binding $I_i$ and the
original provenance established in Phase~2.

For batch $b$, let
\begin{equation}
\mathcal{P}_b=
\{pos_i:VR_i\in\mathcal{B}_{e,b}\}
\end{equation}
denote the authorized executable positions, and let
$\mathcal{L}_b$ and $\mathcal{L}_b'$ denote their old and replacement
leaves, respectively.

\medskip
\noindent\textbf{Step 3: Shared-Frontier Merkle Redaction.}
Rather than verifying and updating one authentication path per
redaction, SFMR derives the minimal external frontier introduced in
Phase~2:
\begin{equation}
\mathcal{F}_b
=
\mathcal{F}_b(\mathcal{P}_b).
\end{equation}
The frontier contains only unaffected sibling nodes required to
reconstruct the paths from all modified leaves to $MR_b$. The PBN first
verifies
\begin{equation}
\mathsf{SFMR.Verify}
(MR_b,\mathcal{P}_b,\mathcal{L}_b,\mathcal{F}_b)=1
\end{equation}
and then computes
\begin{equation}
MR_b'
=
\mathsf{SFMR.Update}
(\mathcal{P}_b,\mathcal{L}_b',
\mathcal{F}_b).
\end{equation}
Internal nodes shared by multiple modified paths are recomputed once.
Hence, the authentication material is proportional to the unique
frontier induced by the selected positions rather than to the sum of
independent Merkle paths. In particular, for $q_b=|\mathcal{P}_b|$,
the independent-path upper cost $O(q_b\log N_b)$ is replaced by work
proportional to the size of the induced authentication subtree and its
frontier, with the largest savings when CORA selects structurally
overlapping positions.

To bind the resulting root to the authorized execution, define
\begin{equation}
\begin{aligned}
\Theta_b
=
H_V\big(&\texttt{SFMR-TRANS}\parallel BID_e\parallel b
\parallel v_b\\
&\parallel MR_b\parallel MR_b'
\parallel \mathsf{enc}(\mathcal{P}_b)
\parallel \Gamma_e\big).
\end{aligned}
\end{equation}

\medskip
\noindent\textbf{Step 4: Authorization-Gated PQCH Adaptation.}
Before releasing a PQCH adaptation share, each participating committee
member $C_k$ independently verifies $Auth_e^B$, $\Gamma_e$, the
authorized membership of every request affecting batch $b$, and the
SFMR transition committed by $\Theta_b$. A share is generated only if
$\mathcal{P}_b$ is a subset of the positions authorized by
$\Gamma_e$ and $MR_b'$ is the deterministic SFMR result for the
corresponding authorized RADC transitions.

Member $C_k$ then computes
\begin{equation}
\delta_{b,k}
\leftarrow
\mathsf{PQCH.PartAdapt}
(td_{e,k},MR_b,r_b,MR_b',\Theta_b).
\end{equation}
After validating at least $t$ distinct shares, the system combines
\begin{equation}
r_b'
\leftarrow
\mathsf{PQCH.Combine}
(\{\delta_{b,k}\}_{k\in\mathcal{T}_b}),
\qquad
|\mathcal{T}_b|\geq t,
\end{equation}
and accepts the adaptation only if
\begin{equation}
\begin{aligned}
\mathsf{PQCH.Hash}(pk_{\mathrm{CH}},MR_b,r_b)
&=
\mathsf{PQCH.Hash}(pk_{\mathrm{CH}},MR_b',r_b')\\
&=CH_b.
\end{aligned}
\end{equation}

Thus, all authorized modifications within batch $b$ produce one SFMR
root transition and one distributed PQCH adaptation. Distinct affected
batches can be prepared in parallel.

\medskip
\noindent\textbf{Step 5: Atomic Authenticated State Finalization.}
For every successfully prepared batch transition,
\begin{equation}
v_b'=v_b+1.
\end{equation}
Each executed transaction updates its SA-RLI entry to
\begin{equation}
E_i'
=
(\tau_i,Tag_i,b,pos_i,PID_i,e,I_i,D_i',ptr_i'),
\end{equation}
where $ptr_i'$ resolves the current authenticated state under $MR_b'$.
Unmodified transaction states remain unchanged.

Let $\Omega_S$ denote the SA-RLI shards containing modified entries.
Only these shards are recomputed:
\begin{equation}
R_s^{*}=
\begin{cases}
R_s', & s\in\Omega_S,\\
R_s,  & s\notin\Omega_S,
\end{cases}
\end{equation}
yielding
\begin{equation}
R_{\mathrm{RLI}}'
=
\mathsf{Merkle.Root}(R_1^{*},\ldots,R_S^{*}),
\qquad
v_{\mathrm{RLI}}'=v_{\mathrm{RLI}}+1.
\end{equation}

For each affected batch, the checkpoint authority constructs
\begin{equation}
\begin{aligned}
A_b'=\big(&b,CH_b,MR_b',r_b',
R_{\mathrm{RLI}}',v_{\mathrm{RLI}}',\\
&e,v_b',ts_b',\sigma_b'\big),
\end{aligned}
\end{equation}
where
\begin{equation}
\begin{aligned}
\sigma_b'=
\mathsf{PQSIG.Sign}\big(
sk_{\mathrm{CP}},
H(&b\parallel CH_b\parallel MR_b'\parallel r_b'\\
&\parallel R_{\mathrm{RLI}}'\parallel v_{\mathrm{RLI}}'
\parallel e\\
&\parallel v_b'\parallel ts_b')\big).
\end{aligned}
\end{equation}
All affected checkpoints reference the same finalized SA-RLI snapshot.
The PBN atomically commits the prepared batch and index transitions
only if their base states remain current; otherwise, the prepared
transitions are discarded and no new finalized state is published.

\medskip
\noindent\textbf{Step 6: Transition-Bound PBRP Generation.}
For each successfully executed request, define the shared batch
execution commitment
\begin{equation}
\begin{aligned}
\Xi_b=
H_V\big(&\texttt{EXEC-EVIDENCE}\parallel
\Theta_b\parallel CH_b\\
&\parallel v_b'\parallel
R_{\mathrm{RLI}}'\parallel v_{\mathrm{RLI}}'\big).
\end{aligned}
\end{equation}
The final Policy-Bound Batch Redaction Proof is
\begin{equation}
\begin{aligned}
PBRP_i=\big(
&PBRP_i^{auth},I_i,D_i,D_i',
b,pos_i,v_b,v_b',\\
&\Theta_b,\Xi_b,A_b',\Pi_i^{\mathrm{SF}}\big),
\end{aligned}
\end{equation}
where $\Pi_i^{\mathrm{SF}}$ is the request-specific membership evidence
linking $(pos_i,L_i,L_i')$ to the shared SFMR transition. The frontier,
PQCH adaptation evidence, and finalized checkpoint are batch-shared
objects and need not be duplicated across PBRPs.

The corresponding compact redaction record is
\begin{equation}
\begin{aligned}
RR_i=\big(
&RID_i,TID_i,BID_e,PID_i,D_i,D_i',\\
&v_b,v_b',H_V(\texttt{PBRP}\parallel
\mathsf{enc}(PBRP_i)),ts_i^{red}\big).
\end{aligned}
\end{equation}
Only successfully finalized transitions produce $RR_i$. The complete
PBRP and its referenced shared evidence are retained for Phase~6
verification and auditing.


\subsubsection*{\textbf{Phase 6: SAVER-Based State-Aware
Verifiable Redaction Auditing}}

This phase introduces \emph{State-Aware Verification Evidence Reuse
(SAVER)}, which organizes finalized PBRP evidence as an authenticated
dependency DAG. Request, authorization, transition, and checkpoint
evidence shared by multiple redactions is represented once and reused
only under the same authenticated state context. Challenged requests
may selectively escalate to independent PQZK verification without
revalidating unaffected shared dependencies.

\medskip
\noindent\textbf{Step 1: Authenticated Redaction Audit Index.}
For each finalized redaction record $RR_i$, the audit service derives
\begin{equation}
\tau_i^A=\mathsf{PRF}_{K_A}(TID_i)
\end{equation}
and constructs
\begin{equation}
E_i^A=
(\tau_i^A,RID_i,BID_e,PID_i,e,v_b,v_b',
h_i^{PBRP},ts_i^{red},ptr_i^A),
\end{equation}
where
\begin{equation}
h_i^{PBRP}
=
H_V(\texttt{PBRP}\parallel\mathsf{enc}(PBRP_i)),
\end{equation}
and $ptr_i^A$ references its supporting evidence. RAI is append-only
with respect to finalized redactions.

Entries are partitioned into $S_A$ authenticated shards:
\begin{equation}
sid_i^A=H_A(\tau_i^A)\bmod S_A,
\end{equation}
whose roots form
\begin{equation}
R_{\mathrm{RAI}}
=
\mathsf{Merkle.Root}(R_1^A,\ldots,R_{S_A}^A).
\end{equation}
A monotonically increasing $v_A$ identifies each snapshot. For audit
epoch $e_A$, the service constructs
\begin{equation}
CP_{e_A}^A
=
H_V(
\texttt{RAI-CP}\parallel e_A\parallel
R_{\mathrm{RAI}}\parallel v_A\parallel ts_A),
\end{equation}
which is anchored to the PBN. Here, $e_A$ is distinct from the
authorization epoch $e$.

\medskip
\noindent\textbf{Step 2: PQ-Authenticated Audit Query Resolution.}
An authorized auditor $AU_j$ submits
\begin{equation}
Q_j^A=
(QID_j,qtype_j,\phi_j,scope_j,e_A,[t_s,t_e],n_j,ts_j)
\end{equation}
and
\begin{equation}
\sigma_j^A
\leftarrow
\mathsf{PQSIG.Sign}(sk_{AU_j},H(Q_j^A)).
\end{equation}
After verifying the auditor, signature, freshness, and scope, the
service resolves
\begin{equation}
\mathcal{R}_{Q_j}
=
\{E_i^A:E_i^A\models Q_j^A\}
\end{equation}
against the selected RAI snapshot. Completeness is claimed only when
the index supports authenticated coverage of the requested predicate
and scope; otherwise, only the returned records are authenticated.

\medskip
\noindent\textbf{Step 3: Authenticated Verification-Dependency DAG.}
For the selected records, SAVER constructs
\begin{equation}
G_{Q_j}^{V}=(V_{Q_j},E_{Q_j}),
\qquad
V_{Q_j}
=
V_{\mathrm{req}}\cup V_{\mathrm{auth}}
\cup V_{\mathrm{trans}}\cup V_{\mathrm{cp}},
\end{equation}
representing request validation, committee authorization, SFMR/PQCH
transitions, and finalized checkpoints, respectively.

Each node $v$ is content-addressed as
\begin{equation}
id(v)=
H_V(
\mathsf{type}(v)\parallel
\mathsf{enc}(\mathsf{payload}(v))\parallel
\mathsf{enc}(\mathsf{parents}(v))).
\end{equation}
Request nodes contain $C_i^{VR}$, $\eta_i$, and
$PBRP_i^{auth}$; authorization nodes contain $Auth_e^B$ and
$\Gamma_e$; transition nodes contain $\Theta_b$ and $\Xi_b$; and
checkpoint nodes contain the corresponding $A_b'$ and RAI checkpoint.
Shared evidence therefore appears as a single dependency node.

The service additionally constructs
\begin{equation}
MP_{Q_j}^A
\leftarrow
\mathsf{Merkle.MultiProof}
(\mathcal{R}_{Q_j},R_{\mathrm{RAI}}),
\end{equation}
together with authenticated coverage evidence $\Gamma_{Q_j}^A$ when
required and supported. The response is
\begin{equation}
Resp_j^A=(Body_j^A,\sigma_{\mathrm{Resp}}^A),
\end{equation}
where
\begin{equation}
\begin{aligned}
Body_j^A=\big(
&QID_j,H(Q_j^A),Ans_j,\mathcal{R}_{Q_j},MP_{Q_j}^A,\\
&G_{Q_j}^{V},\Gamma_{Q_j}^{A},
R_{\mathrm{RAI}},v_A,CP_{e_A}^A\big)
\end{aligned}
\end{equation}
and
\begin{equation}
\sigma_{\mathrm{Resp}}^A
\leftarrow
\mathsf{PQSIG.Sign}
(sk_A,H(\mathsf{enc}(Body_j^A))).
\end{equation}

\medskip
\noindent\textbf{Step 4: State-Bound Verification Reuse.}
The auditor first verifies the response signature, query binding,
anchored RAI checkpoint, multiproof, and any required coverage
evidence, and then traverses $G_{Q_j}^{V}$ in dependency order.

Each node is associated with a typed state context
\begin{equation}
\kappa(v)=
\begin{cases}
(RID_i,h_i^{\mathrm{st}},e),
&v\in V_{\mathrm{req}},\\
(BID_e,\Gamma_e,e),
&v\in V_{\mathrm{auth}},\\
(b,v_b,v_b',\Theta_b),
&v\in V_{\mathrm{trans}},\\
(v_{\mathrm{RLI}}',v_A,CP_{e_A}^A),
&v\in V_{\mathrm{cp}},
\end{cases}
\end{equation}
and verification identifier
\begin{equation}
\chi(v)=
H_V(
\texttt{VERIFY-CTX}\parallel
id(v)\parallel\mathsf{enc}(\kappa(v))).
\end{equation}
After successful verification, the auditor records
\begin{equation}
\mathsf{VCS}[\chi(v)]=\mathsf{valid},
\end{equation}
where VCS denotes the \emph{Verified-Context Store}. Reuse is permitted
only when
\begin{equation}
\mathsf{Reuse}(v,\chi')
=1
\Longleftrightarrow
\chi'=\chi(v)
\land
\mathsf{VCS}[\chi(v)]=\mathsf{valid}.
\label{eq:saver_reuse}
\end{equation}
Thus, any change to the evidence, dependencies, authorization context,
transition state, or checkpoint requires fresh verification.

For each distinct authorization node, the auditor verifies $C_e^B$,
$\Gamma_e$, and at least $t$ committee approvals. For each transition
node, it verifies
\begin{equation}
\mathcal{P}_b^{exec}\subseteq\mathcal{P}_b^{auth},
\end{equation}
the bindings to $\Theta_b$ and $\Xi_b$, and the finalized checkpoint.
Request-specific verification checks
\begin{equation}
I_i'=I_i,\qquad
L_i=H(I_i\parallel D_i),\qquad
L_i'=H(I_i\parallel D_i'),
\end{equation}
together with $\Pi_i^{SF}$ linking
$(pos_i,L_i,L_i')$ to the shared transition. PQCH preservation is
verified once per distinct batch transition:
\begin{equation}
\begin{aligned}
\mathsf{PQCH.Hash}(pk_{\mathrm{CH}},MR_b,r_b)
&=
\mathsf{PQCH.Hash}(pk_{\mathrm{CH}},MR_b',r_b')\\
&=CH_b.
\end{aligned}
\end{equation}

\medskip
\noindent\textbf{Step 5: Dependency-Selective Verification Escalation.}
Routine verification of request $i$ uses
\begin{equation}
\mathcal{D}_i^{\mathrm{routine}}
=
\{
v_i^{\mathrm{req}},
v_e^{\mathrm{auth}},
v_b^{\mathrm{trans}},
v_b^{\mathrm{cp}},
v_A^{\mathrm{cp}}
\}.
\end{equation}
At the request node, the auditor verifies
\begin{equation}
\mathsf{PQSIG.Verify}
\left(
pk_V,
H(RID_i\parallel C_i^{VR}\parallel e),
\alpha_i
\right)=1
\end{equation}
and confirms $h_i^{PQ}=H(\pi_i^{PQ})$. This authenticates the VPS
validation decision but does not independently establish PQZK validity.

For a challenged or forensic request, SAVER expands only the
corresponding request dependency:
\begin{equation}
\mathcal{D}_i^{\mathrm{deep}}
=
\mathcal{D}_i^{\mathrm{routine}}
\cup
\{x_i,\pi_i^{PQ}\},
\end{equation}
and independently verifies
\begin{equation}
\mathsf{PQZK.Verify}
(vk_{\mathrm{PQZK}},x_i,\pi_i^{PQ})=1,
\qquad
H(\pi_i^{PQ})=h_i^{PQ}.
\end{equation}
Unchanged shared dependencies remain reusable according to
Eq.~\eqref{eq:saver_reuse}.

\medskip
\noindent\textbf{Step 6: Query-Specific Audit Decision.}
Finally,
\begin{equation}
\mathsf{AuditVerify}
(PP,Q_j^A,Resp_j^A)
\rightarrow
(ans_j,\beta_j,\ell_j),
\end{equation}
where $\beta_j\in\{0,1\}$ indicates whether the verification conditions
required by $qtype_j$ are satisfied and
$\ell_j\in\{\mathsf{routine},\mathsf{deep}\}$ records the applied
verification level. For exhaustive queries, $\beta_j=1$ additionally
requires authenticated coverage; otherwise, only the authenticity of
the returned records is certified.

For $q$ returned redactions spanning $a$ authorization nodes, $z$
transition nodes, and $c$ checkpoint contexts, with $d$ requests
requiring deep verification, the resulting verification cost is
\begin{equation}
\begin{aligned}
C_{\mathrm{SAVER}}
\approx{}&
qC_{\mathrm{req}}
+aC_{\mathrm{auth}}
+zC_{\mathrm{trans}}\\
&+cC_{\mathrm{cp}}
+dC_{\mathrm{PQZK}},
\end{aligned}
\end{equation}
where typically $a,z,c\ll q$ when redactions share authorization and
execution evidence. SAVER thereby converts the sharing established by
CORA and SFMR into state-bound verification reuse while retaining
request-level escalation when independent policy verification is
required.

    
    
    % =====================================================================
    

\section{Security Analysis}
\label{sec:security}

We analyze VeRedact-PQ against the threat model of Section~III-B.
Let $\mathcal{A}$ be a quantum-polynomial-time (QPT) adversary.
The analysis focuses on the security-critical redaction path; compromise
of the underlying permissioned consensus beyond its prescribed fault
threshold, VPS compromise, endpoint compromise, denial-of-service, and
implementation-level side channels are outside the primary model.

\subsection{Security Model and Assumptions}

VeRedact-PQ relies on the following assumptions.

\begin{itemize}

\item \textbf{A1---Post-Quantum Authentication.}
PQSIG is existentially unforgeable under quantum chosen-message
attacks (EUF-qCMA).

\item \textbf{A2---Private Policy Verification.}
PQZK satisfies quantum computational soundness and zero knowledge for
relation $\mathcal{R}_P$.

\item \textbf{A3---Threshold Redaction.}
PQCH provides collision resistance and threshold-controlled adaptation
against QPT adversaries. Fewer than $t$ contemporaneously valid
trapdoor shares are compromised within each resharing interval.
Previous shares are securely erased after proactive resharing.

\item \textbf{A4---Authenticated Commitments.}
The domain-separated hash functions are collision resistant at the
required post-quantum security level, the PRFs are quantum secure, and
the authenticated Merkle structures are binding.

\item \textbf{A5---Committee and Ledger Thresholds.}
Fewer than $t$ members of the active committee $\mathcal{C}_e$ are
compromised, and the PBN satisfies its prescribed consensus
fault-tolerance condition.

\item \textbf{A6---Validation Service.}
The VPS correctly authenticates requesters, evaluates public policy
conditions, and verifies PQZK proofs before issuing $\alpha_i$.
Routine auditing therefore authenticates the VPS decision rather than
independently reproducing it; deep auditing can independently verify
the committed PQZK proof.

\end{itemize}

Under these assumptions, we establish request admission security,
threshold authorization, authorization-preserving aggregation,
authorized state-transition integrity, end-to-end accountability,
state-bound audit reuse, and the post-quantum security boundary.

\subsection{State-Bound Request Admission}

\textbf{Theorem 1 (Authenticated and State-Bound Admission).}
An adversary cannot obtain a valid Phase~3 validation attestation for
an unauthorized requester, policy, operation, or authenticated
transaction state except with negligible probability.

\emph{Proof:}
A request $R_i$ is admitted only after verifying
\begin{equation}
\mathsf{PQSIG.Verify}
(pk_r,H(R_i),\sigma_i^R)=1.
\label{eq:sec_req_auth}
\end{equation}
Producing a valid signature for a new request without $sk_r$ violates
A1.

The applicable policy is bound by
\begin{equation}
C_{P_i}=H(PID_i\parallel P_i\parallel v_i\parallel e_P),
\end{equation}
while the authenticated transaction state is summarized by
\begin{equation}
h_i^{st}
=
H(TID_i\parallel I_i\parallel D_i\parallel pos_i
\parallel\mathsf{enc}(\mathsf{State}_b)).
\label{eq:sec_state_digest}
\end{equation}
Changing the committed policy or authenticated state while preserving
these bindings requires violating A4.

Private eligibility is verified using
\begin{equation}
\mathsf{PQZK.Verify}
(vk_{\mathrm{PQZK}},x_i,\pi_i^{PQ})=1,
\end{equation}
where
\begin{equation}
x_i=
(ID_r,TID_i,op_i,D_i',C_{P_i},h_i^{st},
e_i,e,ts_r,n_r).
\end{equation}
An adversary without a witness satisfying
$\mathcal{R}_P(x_i,w_i)=1$ can produce an accepting proof only by
violating A2.

Finally,
\begin{equation}
C_i^{VR}
=
H(RID_i\parallel H(R_i)\parallel h_i^{st}
\parallel C_{P_i}\parallel h_i^{PQ}\parallel e),
\end{equation}
where $h_i^{PQ}=H(\pi_i^{PQ})$, is authenticated by
\begin{equation}
\alpha_i=
\mathsf{PQSIG.Sign}
(sk_V,H(RID_i\parallel C_i^{VR}\parallel e)).
\end{equation}
Under A6, this attestation is issued only after the preceding checks.
Thus, successful unauthorized admission requires breaking PQSIG,
PQZK soundness, or an authenticated commitment, each having negligible
probability.

\subsection{Threshold Authorization and CORA Preservation}

\textbf{Theorem 2 (Threshold Authorization Security).}
A coalition controlling fewer than $t$ members of
$\mathcal{C}_e$ cannot independently authorize and execute a valid
redaction except with negligible probability.

\emph{Proof:}
Phase~4 accepts authorization only when at least $t$ committee
approvals authenticate the same
\begin{equation}
C_e^B=
H_V(\texttt{AUTH-BATCH}\parallel BID_e\parallel e\parallel m
\parallel R_e^{VR}\parallel\Gamma_e\parallel ts_e).
\label{eq:sec_auth_batch}
\end{equation}
A coalition of $c<t$ compromised members can produce at most $c$
approvals without forging an uncompromised member's PQ signature,
contradicting A1.

Execution additionally requires at least $t$ valid adaptation shares
\begin{equation}
r_b'=
\mathsf{PQCH.Combine}
(\{\delta_{b,k}\}_{k\in T_b}),
\qquad |T_b|\geq t.
\end{equation}
Under A3, fewer than $t$ contemporaneously valid shares cannot produce
an unauthorized adaptation. Proactive resharing and secure erasure
prevent shares collected across resharing intervals from being
accumulated into a valid threshold set.

\medskip
\noindent\textbf{Lemma 1 (CORA Authorization Preservation).}
CORA cannot make an unvalidated request eligible for committee
authorization or enlarge the authority of a validated request.

\emph{Proof:}
CORA selects exclusively from the validated queue $\mathcal{Q}_e$ and
optimizes only which already validated requests are processed together.
Before authorization, each selected request is revalidated against its
current authenticated state:
\begin{equation}
\widehat h_i^{st}\stackrel{?}{=}h_i^{st}.
\end{equation}
A stale request is therefore excluded from the authorization set.

For the remaining requests,
\begin{equation}
\eta_i=
H_V(\texttt{AUTH-LEAF}\parallel RID_i
\parallel C_i^{VR}\parallel\alpha_i)
\end{equation}
is committed under $R_e^{VR}$, while
\begin{equation}
\Gamma_e=
H_V(\texttt{CORA-PLAN}\parallel
\mathsf{enc}(\{(b_i,pos_i)\}_{i=1}^{m}))
\end{equation}
commits to the authorized execution plan. Both are included in
$C_e^B$. Hence, CORA's cost and overlap optimization affects
\emph{selection} but not request validity or authorization scope.
Insertion, replacement, or relocation after authorization changes
$R_e^{VR}$ or $\Gamma_e$ and invalidates the committee-authorized
commitment under A1 and A4.

\subsection{Authorized-Subset State Transition}

\textbf{Theorem 3 (SFMR Authorized-Subset Integrity).}
Except with negligible probability, a finalized Phase~5 transition
cannot modify a position outside the committee-authorized execution
plan.

\emph{Proof:}
For each affected batch $b$, let
$\mathcal{P}_b^{auth}$ denote the positions authorized by
$\Gamma_e$ and $\mathcal{P}_b^{exec}$ the positions actually executed.
Phase~5 requires
\begin{equation}
\mathcal{P}_b^{exec}\subseteq\mathcal{P}_b^{auth}.
\label{eq:sec_subset}
\end{equation}
Thus, a stale authorized request may be withheld, but an unauthorized
position cannot be introduced.

For each executed request, RADC preserves
\begin{equation}
I_i'=I_i,
\qquad
L_i=H(I_i\parallel D_i),
\qquad
L_i'=H(I_i\parallel D_i').
\end{equation}
Changing the transaction identifier, owner, data type, policy,
timestamp, or tag while preserving $I_i$ violates A4.

SFMR verifies the old leaves and external authentication frontier
against $MR_b$, recomputes the affected shared subtree, and obtains
$MR_b'$. The transition is committed by
\begin{equation}
\begin{aligned}
\Theta_b=H_V(&\texttt{SFMR-TRANS}\parallel BID_e\parallel b
\parallel v_b\parallel MR_b\parallel MR_b'\\
&\parallel\mathsf{enc}(\mathcal{P}_b^{exec})
\parallel\Gamma_e).
\end{aligned}
\label{eq:sec_theta}
\end{equation}
Consequently, substituting an executed position, old root, new root,
or authorization plan changes $\Theta_b$ unless A4 is violated.

Committee members generate $\delta_{b,k}$ only after verifying the
authorization, $\Gamma_e$, the subset relation in
Eq.~\eqref{eq:sec_subset}, and the SFMR transition. The finalized
adaptation must satisfy
\begin{equation}
\begin{split}
\mathsf{PQCH.Hash}(pk_{\mathrm{CH}},MR_b,r_b)
={}&
\mathsf{PQCH.Hash}(pk_{\mathrm{CH}},MR_b',r_b')\\
={}&CH_b.
\end{split}
\label{eq:sec_pqch_continuity}
\end{equation}
An unauthorized transition therefore requires violating the Merkle
binding, $\Gamma_e/\Theta_b$ binding, threshold authorization, or
PQCH security. Under A1--A5, each event has negligible probability.

\subsection{Freshness and Replay Resistance}

\textbf{Theorem 4 (State and Epoch Freshness).}
Validation or authorization evidence generated for one authenticated
state or authorization epoch cannot be reused for a different state or
expired epoch except with negligible probability.

\emph{Proof:}
Phase~3 binds the request to $h_i^{st}$ and authorization epoch $e$.
Before committee authorization, the current state is recomputed and
must satisfy
\begin{equation}
\widehat h_i^{st}=h_i^{st}.
\end{equation}
After successful execution,
\begin{equation}
v_b'=v_b+1,
\qquad
v_{\mathrm{RLI}}'=v_{\mathrm{RLI}}+1,
\end{equation}
and the new checkpoint commits to the resulting roots and versions.

Therefore, a request validated against an earlier state produces a
different state digest after the relevant state changes. Likewise,
$\alpha_i$, $C_e^B$, and $\Gamma_e$ are bound to the authorization
epoch, whose validity is checked before execution. Rewriting these
values while preserving existing authentication requires a hash
collision or PQ-signature forgery. Hence stale or cross-epoch evidence
is rejected except with negligible probability.

\subsection{End-to-End Redaction Accountability}

\textbf{Theorem 5 (Policy-Bound Redaction Accountability).}
Every finalized VeRedact-PQ redaction can be cryptographically linked
to its validated request, threshold authorization, authorized execution
plan, executed state transition, and finalized checkpoint.

\emph{Proof:}
For each executed request, $PBRP_i^{auth}$ proves membership of
$\eta_i$ in the committee-authorized $R_e^{VR}$. The authorization
also commits to $\Gamma_e$, which defines the authorized batch/position
plan.

The executed transition is bound by $\Theta_b$, and its finalized
evidence is summarized as
\begin{equation}
\Xi_b=
H_V(\texttt{EXEC-EVIDENCE}\parallel
\Theta_b\parallel CH_b\parallel v_b'
\parallel R_{\mathrm{RLI}}'\parallel v_{\mathrm{RLI}}').
\end{equation}
Request-specific SFMR evidence $\Pi_i^{SF}$ links
$(pos_i,L_i,L_i')$ to this shared transition. Hence
\begin{equation}
PBRP_i=
(PBRP_i^{auth},I_i,D_i,D_i',b,pos_i,
v_b,v_b',\Theta_b,\Xi_b,A_b',\Pi_i^{SF})
\end{equation}
forms a cryptographic chain from request authorization to finalized
execution.

The redaction record commits to
\begin{equation}
h_i^{PBRP}
=
H_V(\texttt{PBRP}\parallel\mathsf{enc}(PBRP_i)),
\end{equation}
which is subsequently authenticated by RAI and its blockchain-anchored
checkpoint. Altering any component without changing the corresponding
commitment violates A4, while forging its authenticated authorization
or checkpoint violates A1 or the PBN assumption. Therefore, accepted
redactions remain accountable except with negligible probability.


\subsection{SAVER State-Bound Verification Reuse}

\textbf{Theorem 6 (State-Bound Verification-Reuse Soundness).}
Assume that $\mathsf{VCS}[\chi(v)]=\mathsf{valid}$ is recorded only
after successful verification of evidence node $v$ under context
$\chi(v)$. A valid result cannot be reused for different evidence or a
different authenticated state context except with negligible
probability.

\emph{Proof:}
Each evidence node is content-addressed by
\begin{equation}
id(v)=
H_V(\mathsf{type}(v)\parallel
\mathsf{enc}(\mathsf{payload}(v))\parallel
\mathsf{enc}(\mathsf{parents}(v))),
\end{equation}
and bound to its typed state context by
\begin{equation}
\chi(v)=
H_V(\texttt{VERIFY-CTX}\parallel
id(v)\parallel\mathsf{enc}(\kappa(v))),
\label{eq:sec_saver_ctx}
\end{equation}
where $\kappa(v)$ contains the state variables relevant to the node
type.

A verified result is reusable only if
\begin{equation}
\mathsf{Reuse}(v,\chi')=1
\Longleftrightarrow
\chi'=\chi(v)
\land
\mathsf{VCS}[\chi(v)]=\mathsf{valid}.
\label{eq:sec_reuse}
\end{equation}
Any change to the node payload, parent dependencies, authorization
context, transition state, version, or checkpoint changes either
$id(v)$ or $\kappa(v)$. Reusing the previous result while preserving
the same $\chi(v)$ would therefore require a collision in $H_V$,
contradicting A4. Hence, SAVER reuses verification results only for
identical authenticated evidence and state contexts.

\medskip
\noindent\textbf{Corollary 1 (Dependency-Selective Escalation).}
Independent verification of a challenged request does not require
reverification of unchanged shared dependencies.

For routine auditing, the verifier authenticates the VPS decision by
\begin{equation}
\mathsf{PQSIG.Verify}
(pk_V,H(RID_i\parallel C_i^{VR}\parallel e),\alpha_i)=1
\end{equation}
and confirms
\begin{equation}
H(\pi_i^{PQ})=h_i^{PQ}.
\end{equation}
This establishes authenticated evidence that the trusted VPS accepted
the corresponding state-bound PQZK validation, but does not
independently establish PQZK validity.

For deep auditing, the challenged request branch is expanded and the
verifier independently checks
\begin{equation}
\mathsf{PQZK.Verify}
(vk_{\mathrm{PQZK}},x_i,\pi_i^{PQ})=1,
\qquad
H(\pi_i^{PQ})=h_i^{PQ}.
\end{equation}
Shared authorization, transition, and checkpoint dependencies remain
reusable whenever their contexts satisfy
Eq.~\eqref{eq:sec_reuse}. Therefore, selective escalation provides
independent verification of the challenged request without weakening
the soundness of reused shared evidence.

\subsection{Audit Integrity, Completeness, and Privacy}

\textbf{Theorem 7 (Authenticated and Privacy-Preserving Audit).}
Under A1, A2, and A4, an authorized auditor can verify authenticated
redaction evidence without learning the private policy witness, while
modification, substitution, or stale reuse of authenticated audit
evidence is detectable except with negligible probability.

\emph{Proof:}
Private credentials and attributes occur in witness $w_i$. By PQZK
zero knowledge, $\pi_i^{PQ}$ reveals no information about $w_i$ beyond
the validity of $\mathcal{R}_P(x_i,w_i)$. Routine auditing need not
process $w_i$ and can rely on the authenticated VPS attestation;
deep auditing reveals the PQZK proof but not its witness.

RAI addresses transaction entries using
\begin{equation}
\tau_i^A=\mathsf{PRF}_{K_A}(TID_i)
\end{equation}
and authenticates the resulting entries under
$R_{\mathrm{RAI}}$. The audit snapshot is bound to
\begin{equation}
CP_{e_A}^{A}
=
H_V(\texttt{RAI-CP}\parallel e_A
\parallel R_{\mathrm{RAI}}\parallel v_A\parallel ts_A),
\end{equation}
which is anchored to the PBN. The audit response is additionally
PQ-signed. Modifying an entry, proof, root, version, or checkpoint
therefore requires breaking the authenticated Merkle structure,
$H_V$, PQSIG, or the anchored ledger assumption.

This privacy claim is intentionally limited: VeRedact-PQ protects the
private authorization witness and avoids requiring plaintext
transaction content for verification, but does not claim to hide all
query/access patterns or metadata revealed by authorized audit results.

\medskip
\noindent\textbf{Remark on Audit Completeness.}
A Merkle membership or multiproof establishes authenticity of returned
RAI entries but does not by itself prove completeness for an arbitrary
semantic predicate. VeRedact-PQ therefore claims result-set
completeness only when the corresponding RAI query is accompanied by
authenticated coverage, range, or non-membership evidence
$\Gamma_{Q_j}^{A}$. Otherwise, the guarantee is limited to integrity
and freshness of the returned records.

\subsection{Post-Quantum Security Boundary}

\textbf{Theorem 8 (Post-Quantum Redaction Security).}
Under A1--A6 and the stated committee and PBN thresholds, a QPT
adversary cannot cause an unauthorized VeRedact-PQ redaction to be
accepted except with negligible probability.

\emph{Proof:}
An accepted redaction requires the following security-critical chain:
\begin{equation}
\begin{split}
&\mathsf{PQ\ requester\ authentication}
\Rightarrow
\mathsf{state\mbox{-}bound\ PQZK\ validation}\\
&\Rightarrow
t\text{-of-}n\ \mathsf{PQ\ committee\ authorization}
\Rightarrow
\mathsf{authorized\ SFMR\ transition}\\
&\Rightarrow
\mathsf{threshold\ PQCH\ adaptation}
\Rightarrow
\mathsf{authenticated\ PQ\ audit\ evidence}.
\end{split}
\label{eq:sec_pq_chain}
\end{equation}
Theorems~1--7 show that bypassing this chain requires violating at
least one of PQSIG unforgeability, PQZK soundness, hash/Merkle binding,
threshold PQCH security, the committee threshold, or the assumed PBN
consensus condition. Each cryptographic event occurs only with
negligible probability under A1--A5.

The result does not imply that the underlying permissioned blockchain
is itself fully post-quantum secure. Native membership, transport,
endorsement, or consensus credentials may remain classical. Compromise
of such a credential alone does not provide a valid requester
signature, VPS attestation, $t$ committee approvals, or $t$ PQCH
adaptation shares. However, if a quantum adversary compromises enough
native identities to violate the PBN consensus assumption, ledger
integrity or availability may fail independently of VeRedact-PQ.

Therefore, VeRedact-PQ establishes a post-quantum security boundary for
the security-critical redaction and audit path, subject to the stated
committee, VPS, and permissioned-consensus assumptions.

    
    % (Evaluation skeleton in VedRedact.tex)
    % [TBD] = value to be filled from measurements / your decision.
    % "% TODO-VERIFY" = baseline entry inferred from Table I features;
    %                   check against the cited paper's own cost analysis.
    % =====================================================================
    
    \section{Evaluation}
    \label{sec:evaluation}
    
    This section evaluates VeRedact-PQ through analytical computation- and
    communication-cost comparison and experimental performance analysis.
    Four representative redactable blockchain schemes are used as
    baselines: Scheme~\cite{ref1} for decentralized threshold
    redaction with a consistency check; Scheme~\cite{ref13} for auditable redaction
    with version-freshness proofs under a trusted system manager;
    Scheme~\cite{ref27} for robust $(t,n)$-threshold redaction;
    and Scheme~\cite{ref34} for attribute-based,
    policy-hiding redaction. Together, these schemes cover complementary
    capabilities in distributed authorization, policy control, and
    verifiable auditing, as summarized in Table~\ref{tab:comparison}. Each
    baseline enters only the experiments whose stage its paper defines:
    Schemes~\cite{ref13} and~\cite{ref27} define no authorization step
    separate from redaction, and Schemes~\cite{ref27} and~\cite{ref34}
    define no audit protocol.
    % ---------------------------------------------------------------------

\subsection{Computation Cost Analysis}
\label{subsec:computation-cost}

This subsection compares the dominant computation costs of VeRedact-PQ
with the primary baselines across request validation, authorization,
redaction execution, and audit verification. Consider a CORA-selected
set of $m$ validated requests affecting $|\Omega_e|$ transaction
batches, each containing $N$ leaves. For each affected batch $b$, let
$f_b=|\mathcal{F}_b(\mathcal{P}_b^{exec})|$ denote the SFMR frontier
size, and let $F_e=\sum_{b\in\Omega_e}f_b$. An audit query returns
$n_Q$ redaction records involving $a_Q$ distinct authorization nodes,
$z_Q$ distinct transition nodes, and $c_Q$ distinct checkpoint
contexts. Of these records, $d_Q$ may be escalated to independent PQZK
verification. Table~\ref{tab:cost-notation} summarizes the notation,
and Table~\ref{tab:comp-cost} compares the dominant costs. N/A denotes
a stage not explicitly supported by the corresponding scheme.

\begin{table}[t]
\centering
\caption{Notation Used in the Cost Analysis}
\label{tab:cost-notation}
\scriptsize
\renewcommand{\arraystretch}{1.08}
\setlength{\tabcolsep}{3pt}
\begin{tabular}{ll}
\hline
\textbf{Notation} & \textbf{Description} \\
\hline
$m$ & Requests selected by CORA \\
$|\Omega_e|$ & Transaction batches affected by CORA selection \\
$N$ & Leaves per transaction batch \\
$N_I$, $N_A$ & SA-RLI and RAI entries \\
$n$, $t$ & Committee size and authorization threshold \\
$f_b$, $F_e$ & SFMR frontier size for batch $b$, total frontier size \\
$n_Q$ & Redaction records returned by an audit query \\
$a_Q$ & Distinct authorization dependencies in audit result \\
$z_Q$ & Distinct transition dependencies in audit result \\
$c_Q$ & Distinct checkpoint contexts in audit result \\
$d_Q$ & Requests escalated to deep PQZK verification \\
$c$ & Challenged blocks in~\cite{ref13}, $c\le n_Q$ \\
\hline
$T_H$, $T_{\mathrm{PRF}}$ & Hash and PRF evaluation \\
$T_S$, $T_V$ & PQ signature generation and verification \\
$T_{ZP}$, $T_{ZV}$ & PQZK proof generation and verification \\
$T_{CH}$ & PQCH hash evaluation \\
$T_{PA}$, $T_{CB}$ & PQCH partial adaptation, share combination \\
$T_V^{c}$ & Classical signature verification \\
$T_{AD}^{c}$, $T_{CV}^{c}$ & Classical CH adaptation, verification \\
$T_{PA}^{c}$, $T_{CB}^{c}$ & Classical CH partial adaptation, combination \\
$T_{Pol}$ & Policy-based CH authorization operation \\
$T_{Acc}$ & RSA-accumulator membership/update operation \\
$T_{Tag}$ & Identity-based RSA tag operation~\cite{ref13} \\
$T_E^{c}$ & Classical modular exponentiation \\
\hline
\end{tabular}
\end{table}


\begin{table*}[t]
\centering
\caption{Computation-Cost Comparison of the Redaction Lifecycle}
\label{tab:comp-cost}
\footnotesize
\renewcommand{\arraystretch}{1.25}
\setlength{\tabcolsep}{3.8pt}
\begin{tabular}{lcccc}
\toprule
\textbf{Scheme} &
\makecell{\textbf{Request Validation}\\\textbf{(per request)}} &
\makecell{\textbf{Authorization}\\\textbf{(per $m$ requests)}} &
\makecell{\textbf{Redaction Execution}\\\textbf{(per $m$ requests)}} &
\makecell{\textbf{Audit Verification}\\\textbf{(per $n_Q$ records)}} \\
\midrule

Scheme~\cite{ref1}
& $T_{Acc}+2T_V^{c}$
& $m\,t(T_{Acc}+2T_V^{c})$
& \makecell{$m(tT_{PA}^{c}+T_{CB}^{c}+T_{CV}^{c}+T_{Acc})$\\
$+O(m\log N)T_H$}
& $n_Q(T_{CV}^{c}+T_{Acc})$ \\

Scheme~\cite{ref13}
& N/A
& N/A
& \makecell{$m(T_{AD}^{c}+T_{CV}^{c}+3T_{Tag}+T_{Acc})$\\
$+O(mN)T_H$}
& $O(c)T_E^{c}$ \\

Scheme~\cite{ref27}
& N/A
& N/A
& $m(tT_{PA}^{c}+T_{CB}^{c}+T_{CV}^{c}+T_V^{c})$
& N/A \\

Scheme~\cite{ref34}
& N/A
& $mT_{Pol}$
& $m(T_{AD}^{c}+2T_{CV}^{c})$
& N/A \\

\midrule
\textbf{VeRedact-PQ}
& \makecell{$T_V+T_{ZV}+T_S+T_{\mathrm{PRF}}$\\
$+O(\log N_I+\log N)T_H$}
& \makecell{$mT_V+t(T_S+T_V)$\\
$+O(m)T_H$}
& \makecell{$|\Omega_e|(tT_{PA}+T_{CB}+2T_{CH}+T_S)$\\
$+O(m+F_e+\Delta_I)T_H$}
& \makecell{$(n_Q+a_Qt+c_Q+1)T_V+2z_QT_{CH}$\\
$+d_QT_{ZV}
+O(n_Q+a_Q+z_Q+c_Q+\log N_A)T_H$} \\

\bottomrule
\multicolumn{5}{p{0.96\textwidth}}{\footnotesize
Requester-side proof/signature generation ($T_{ZP}+T_S$) is excluded
from server-side request validation. $F_e$ is the total SFMR frontier
size and $\Delta_I$ denotes the authenticated SA-RLI update work
($O(m\log N_I)$ without shared index-update paths). Routine SAVER
verification has $d_Q=0$; deep auditing adds $d_QT_{ZV}$.
Authorization, transition, and checkpoint evidence is verified once
per distinct SAVER dependency. The additional $T_V$ in audit
verification authenticates the audit-service response. In~\cite{ref1},
a block can be redacted only once and approvals carry no signature;
in~\cite{ref27}, the threshold is enforced inside adaptation but no
separate request-validation or audit stage is defined; the audit
in~\cite{ref13} produces an aggregate decision for the challenged
blocks; one-off attribute-key issuance in~\cite{ref34} is excluded.}
\end{tabular}
\end{table*}

Table~\ref{tab:comp-cost} shows that the baselines incur repeated
cryptographic costs at different stages. Scheme~\cite{ref1} repeats
request checking and threshold adaptation per redaction;
Scheme~\cite{ref13} incurs per-request chameleon-hash, tag,
accumulator, and Merkle costs; Scheme~\cite{ref27} performs threshold
adaptation for every redaction without explicit validation or auditing;
and Scheme~\cite{ref34} applies policy and chameleon-hash operations
per request but provides no audit mechanism.

VeRedact-PQ incurs a higher admission cost because each request
undergoes state-bound PQZK verification once before CORA selection.
Subsequent authorization uses the compact VPS attestation, while the
$t$ committee approvals are shared across the selected set:
\begin{equation}
C_{\mathrm{auth}}
=
mT_V+t(T_S+T_V)+O(m)T_H.
\end{equation}
Thus, the threshold component amortizes to
$t(T_S+T_V)/m$ per request.

CORA and SFMR reduce execution cost by exploiting batch and Merkle-path
locality. SFMR processes shared authentication structure according to
the total frontier size $F_e$, while threshold PQCH adaptation is
performed once per affected batch:
\begin{equation}
C_{\mathrm{PQCH}}
=
|\Omega_e|(tT_{PA}+T_{CB}+2T_{CH}),
\qquad |\Omega_e|\leq m.
\end{equation}
Hence, clustered requests reduce both $|\Omega_e|$ and $F_e$ relative
to independent processing.

SAVER similarly avoids repeated audit verification by sharing
state-bound dependencies:
\begin{equation}
C_{\mathrm{SAVER}}
=
n_QC_{\mathrm{req}}
+a_QC_{\mathrm{auth}}
+z_QC_{\mathrm{trans}}
+c_QC_{\mathrm{cp}}
+d_QT_{ZV},
\label{eq:saver_cost}
\end{equation}
where typically $a_Q,z_Q,c_Q\ll n_Q$. Routine auditing has $d_Q=0$;
only challenged requests require independent PQZK verification.

Overall, VeRedact-PQ does not assume that PQ primitives are cheaper
than classical ones. Its scalability instead comes from reducing their
invocation frequency: PQZK verification occurs once per admitted
request, committee authorization once per CORA-selected set, PQCH
adaptation once per affected batch, and shared audit evidence once per
matching SAVER state context. The experiments below quantify these
amortization benefits at approximately 128-bit security.


\subsection{Communication and Storage Cost Analysis}
\label{subsec:communication-cost}

Let $|R|$, $|\sigma|$, $|\pi|$, and $|h|$ denote the sizes of a
redaction request, PQ signature, PQZK proof, and hash digest,
respectively. Let $|\sigma^{c}|$, $|r^{c}|$, $|ID|$, $|BH|$, and
$|\mathbb{Z}_N|$ denote a classical signature, classical
chameleon-hash randomness, node identity, block header, and RSA-group
element. We use the same workload notation as in
Section~\ref{subsec:computation-cost}. Table~\ref{tab:comm-cost}
compares the dominant communication and persistent-state costs;
redacted payloads common to all schemes are omitted.

\begin{table*}[t]
\centering
\caption{Communication and Storage Cost Comparison}
\label{tab:comm-cost}
\footnotesize
\renewcommand{\arraystretch}{1.25}
\setlength{\tabcolsep}{3.8pt}
\begin{tabular}{lcccc}
\toprule
\textbf{Scheme} &
\makecell{\textbf{Request Submission}\\\textbf{(per request)}} &
\makecell{\textbf{Authorization Evidence}\\\textbf{(per $m$ requests)}} &
\makecell{\textbf{On-Chain Redaction State}\\\textbf{(per $m$ requests)}} &
\makecell{\textbf{Audit Response}\\\textbf{(per $n_Q$ records)}} \\
\midrule

Scheme~\cite{ref1}
& $|R|+|\sigma^{c}|$
& $mt|ID|$
& $m(O(1)|h|+|r^{c}|)$
& $n_Q(|BH|+|\mathbb{Z}_N|)$ \\

Scheme~\cite{ref13}
& $|R|$
& N/A
& $m(|BH|+3|\mathbb{Z}_N|)$
& $c|BH|+(c+10)|\mathbb{Z}_N|$ \\

Scheme~\cite{ref27}
& N/A
& N/A
& $m|r^{c}|$
& N/A \\

Scheme~\cite{ref34}
& N/A
& N/A
& $m|r^{c}|$
& N/A \\

\midrule
\textbf{VeRedact-PQ}
& $|R|+|\sigma|+|\pi|$
& \makecell{$t|\sigma|+O(m)|h|$\\
$+O(F_e)|h|$}
& $|\Omega_e|(|\sigma|+O(1)|h|)$
& \makecell{$(n_Q+a_Qt+c_Q)|\sigma|$\\
$+O(n_Q+a_Q+z_Q+c_Q+\log N_A)|h|$\\
$+d_Q|\pi|$} \\

\bottomrule
\multicolumn{5}{p{0.96\textwidth}}{\footnotesize
For VeRedact-PQ, complete PQZK proofs, PBRP objects, and detailed SFMR
evidence may be retained in the off-ledger evidence store, while their
digests or authenticated roots are anchored on-chain. Routine SAVER
auditing has $d_Q=0$; deep auditing transfers the PQZK proofs only for
the $d_Q$ challenged requests. Redacted payloads common to all schemes
are omitted. The $|QR_N|$ elements in~\cite{ref13} are represented by
$|\mathbb{Z}_N|$.}
\end{tabular}
\end{table*}

Table~\ref{tab:comm-cost} shows that VeRedact-PQ incurs larger
per-request cryptographic objects because request submission includes a
PQ signature and PQZK proof. This cost is incurred once during
validation. Subsequent phases reference the validated request through
compact commitments and the VPS attestation rather than retransmitting
the complete PQZK proof. Committee authorization further contributes
$t|\sigma|$ once per CORA-selected set instead of $t$ approvals per
request. The authorization root $R_e^{VR}$ and CORA plan $\Gamma_e$
bind the selected requests and execution locality, while request-level
membership evidence can be retained with PBRP rather than replicated
on-chain.

CORA and SFMR also reduce execution-related communication. For
independent processing, $m$ requests may carry overlapping Merkle
authentication paths and require up to $m$ separate chameleon-hash
adaptations. In VeRedact-PQ, shared authentication structure is
represented by the SFMR frontier, whose size is summarized by $F_e$,
and only one threshold PQCH adaptation is required per affected
transaction batch. Consequently, persistent on-chain transition state
scales primarily with
\begin{equation}
|\Omega_e|,
\qquad |\Omega_e|\leq m,
\end{equation}
because one finalized checkpoint $A_b'$ is anchored per affected batch.
Detailed request-level PBRP evidence can remain off-ledger and is
integrity-bound through authenticated digests and RAI checkpoints.

SAVER extends this sharing to audit communication. Without dependency
reuse, an audit over $n_Q$ records may repeatedly transmit the same
committee approvals, transition evidence, and checkpoint information.
SAVER instead represents common authorization, transition, and
checkpoint objects once for each distinct dependency. Its dominant
response size can therefore be expressed as
\begin{equation}
\begin{aligned}
S_{\mathrm{SAVER}}
={}&O(n_Q)|h|
+a_Q\,O(t|\sigma|)
+z_Q\,O(|h|)\\
&+c_Q\,O(|\sigma|+|h|)
+d_Q|\pi|
+S_{\mathrm{RAI}},
\end{aligned}
\label{eq:saver_comm}
\end{equation}
where $S_{\mathrm{RAI}}$ denotes the query-scoped RAI authentication
evidence. For related redactions, typically
$a_Q,z_Q,c_Q\ll n_Q$, so shared evidence is transferred once rather
than duplicated across records. Routine auditing has $d_Q=0$, whereas
only challenged requests incur the additional PQZK-proof transfer.

The persistent storage design follows the same principle. SA-RLI and
RAI require $O(N_I)$ and $O(N_A)$ authenticated index storage,
respectively, while the blockchain retains compact roots, versions,
authorization commitments, and finalized checkpoints. Large PQZK,
PBRP, and SFMR objects may be stored off-ledger, but their integrity
remains bound to blockchain-anchored commitments. Thus, VeRedact-PQ
does not reduce the size of individual post-quantum objects; rather, it
limits their replication by storing and communicating shared
authorization, transition, and verification evidence at the finest
necessary granularity.

    % ---------------------------------------------------------------------

\subsection{Performance Analysis}
\label{subsec:performance-analysis}

This subsection evaluates VeRedact-PQ under varying redaction
workloads, committee configurations, request locality, audit scopes,
and on-chain anchoring costs. The experiments examine whether CORA,
SFMR, and SAVER reduce repeated authorization, execution, and
verification work while preserving the request-level validation and
post-quantum security properties established in
Section~\ref{sec:security}.

\subsubsection{Experimental Setup}

VeRedact-PQ was implemented in approximately \textbf{[TBD]} lines of
\textbf{[TBD: language]} code and evaluated on a server equipped with
\textbf{[TBD: CPU, cores, clock]}, \textbf{[TBD]}\,GB RAM, running
Ubuntu \textbf{[TBD]}. PBN nodes, committee members, the VPS, and
auditors were deployed in Docker containers. Inter-node latency was
configured to \textbf{[TBD]}\,ms using Linux \texttt{tc netem} to
emulate a multi-organization deployment.

PQSIG was instantiated with ML-DSA-65 (FIPS~204) using the Open
Quantum Safe \texttt{liboqs} library. $H$, $H_1$, $H_2$, $H_A$, and
$H_V$ were implemented using domain-separated SHA3-256, while the
SA-RLI and RAI PRFs use HMAC-SHA3-256 with 256-bit keys
$K_{\mathrm{idx}}$ and $K_A$. PQZK was instantiated using a
transparent hash-based STARK system \textbf{[TBD: library]},
consistent with the trapdoor-free setup assumed in Phase~1. PQCH was
implemented using an SIS-based post-quantum chameleon-hash construction
following~\cite{ref10,ref17}; the threshold adaptation mechanism uses
\textbf{[TBD: exact $t$-of-$n$ construction]}. SA-RLI, RAI, and the
SFMR authenticated structures use binary SHA3-256 Merkle trees with
multiproof support.

The PBN was implemented as a \textbf{[TBD]}-validator Hyperledger Besu
network using QBFT consensus. Solidity contracts anchor policy
commitments, committee configurations, finalized batch checkpoints,
and RAI checkpoints. PQ signature verification is performed by
validator-side services before anchoring rather than within the smart
contracts. Table~\ref{tab:primitives} reports the measured primitive
costs used in the evaluation.

The workload contains synthetic enterprise transactions with payloads
from 256\,B to 4\,KB, grouped into batches of $N$ leaves. Redaction
targets follow a Zipf distribution with skew $s$, where $s=0$
approximates uniform placement and increasing $s$ concentrates requests
on fewer transaction batches. This parameter directly controls the
batch and Merkle-path locality available to CORA and SFMR. Request
arrivals follow a Poisson process, supplemented by bursty on/off
workloads to evaluate CORA under changing queue pressure. Audit
workloads vary the number of returned records, shared authorization and
transition dependencies, and the fraction of requests escalated to
deep PQZK verification.

Unless otherwise stated, Table~\ref{tab:params} gives the default
configuration. Each experiment is repeated in three independent runs
with different random seeds. Within a run, each measurement point
collects 30 samples (100 for the primitives in Table~\ref{tab:primitives},
5 for Ref.~\cite{ref13} at $n_Q=10^4$, and every request of the 60-s
measurement window in Experiment~1); each reported value is the mean
over the three runs of the per-run median (p95 likewise), with a 95\%
confidence interval over the runs (Student's $t$, two degrees of
freedom).

To separate architectural effects from cryptographic primitive costs,
the baselines are reimplemented in the same experimental environment
using their prescribed classical constructions at approximately
128-bit security. Operations or lifecycle stages not defined by the
original scheme are reported as unsupported rather than supplied by
VeRedact-PQ. Otherwise, each baseline follows its specified
authorization, redaction, and audit workflow.

\begin{table}[t]
\centering
\caption{Execution Time of Cryptographic Primitives}
\label{tab:primitives}
\scriptsize
\renewcommand{\arraystretch}{1.08}
\setlength{\tabcolsep}{4pt}
\begin{tabular}{llc}
\hline
\textbf{Operation} & \textbf{Instantiation} & \textbf{Time (ms)} \\
\hline
\multicolumn{3}{l}{\textit{VeRedact-PQ}} \\
$T_S$ / $T_V$
& ML-DSA-65 sign / verify
& [TBD] / [TBD] \\
$T_{ZP}$ / $T_{ZV}$
& PQZK prove / verify
& [TBD] / [TBD] \\
$T_{CH}$
& PQCH hash
& [TBD] \\
$T_{PA}$ / $T_{CB}$
& PQCH partial adapt / combine
& [TBD] / [TBD] \\
$T_H$
& SHA3-256 (64-byte input)
& [TBD] \\
$T_{\mathrm{PRF}}$
& HMAC-SHA3-256
& [TBD] \\
\hline
\multicolumn{3}{l}{\textit{Baselines}} \\
$T_V^{c}$
& ECDSA secp256k1 verify~\cite{ref1,ref27}
& [TBD] \\
$T_{PA}^{c}$ / $T_{CV}^{c}$
& Improved DCH, BLS12-381~\cite{ref1}
& [TBD] / [TBD] \\
$T_{PA}^{c}$ / $T_{CV}^{c}$
& ETCH, secp256k1~\cite{ref27}
& [TBD] / [TBD] \\
$T_{AD}^{c}$ / $T_{CV}^{c}$
& Double-trapdoor CH, RSA-3072~\cite{ref13}
& [TBD] / [TBD] \\
$T_{AD}^{c}$ / $T_{CV}^{c}$
& CHET, RSA-3072~\cite{ref34}
& [TBD] / [TBD] \\
$T_{Acc}$
& RSA-3072 accumulator~\cite{ref1,ref13}
& [TBD] \\
$T_{Tag}$
& Identity-based RSA-3072 tag~\cite{ref13}
& [TBD] \\
$T_{Pol}$
& MA-ABE decryption, BLS12-381~\cite{ref34}
& [TBD] \\
\hline
\end{tabular}
\end{table}

\begin{table}[t]
\centering
\caption{Default Experimental Parameters}
\label{tab:params}
\scriptsize
\renewcommand{\arraystretch}{1.08}
\setlength{\tabcolsep}{4pt}
\begin{tabular}{ll}
\hline
\textbf{Parameter} & \textbf{Default value} \\
\hline
Committee size $n$ / threshold $t$ & [7] / [5] \\
CORA bounds $B_{\min}$ / $B_{\max}$ & [8] / [256] \\
Maximum waiting time $T_{\max}$ & [500]\,ms \\
Transaction leaves per batch $N$ & [256] \\
SA-RLI shards $S$ / RAI shards $S_A$ & [16] / [16] \\
Ledger size & [$10^5$] transactions \\
Zipf skew $s$ & [0.8] \\
Deep-audit fraction & [10\%] \\
Signature scheme & ML-DSA-65 \\
Repetitions & 3 runs, 30 samples per point, 95\% CI \\
\hline
\end{tabular}
\end{table}

    
    % ------------------------------------
\subsubsection*{\textbf{Experiment 1: End-to-End Redaction Performance}}

This experiment evaluates end-to-end redaction performance under
increasing request rates and transaction-batch locality. The arrival
rate is varied from \textbf{[100]} to \textbf{[5,000]} requests/s.
End-to-end latency is measured from request submission to finalization
of the corresponding Phase~5 checkpoint, while throughput is measured
as finalized redactions per second. VeRedact-PQ is compared with
Refs.~\cite{ref1}, \cite{ref13}, \cite{ref27}, and~\cite{ref34}.

To isolate the contributions of CORA and SFMR, we additionally evaluate
the following internal variants:
\begin{enumerate}
\item \textbf{Per-Request}: each validated request is authorized and
executed independently, without CORA or SFMR;
\item \textbf{Fixed-Batch}: validated requests are grouped in arrival
order using a fixed batch size of \textbf{[64]}, without CORA's
cost/overlap-aware selection; and
\item \textbf{CORA-No-SFMR}: CORA selects requests as in VeRedact-PQ,
but each authorized redaction performs an independent Merkle update
and PQCH adaptation.
\end{enumerate}
The full \textbf{VeRedact-PQ} configuration combines CORA selection
with SFMR shared-frontier execution.

To examine locality, the Zipf skew is varied as
$s\in\{0,0.4,0.8,1.2\}$ at a fixed arrival rate. We record the number
of affected transaction batches $|\Omega_e|$, PQCH adaptations, and
total SFMR frontier size $F_e$ per 1,000 finalized redactions. These
metrics directly quantify the execution sharing predicted by the cost
analysis.

\begin{figure}[t]
\centering
\includegraphics[width=0.8\linewidth]{exp1_redaction_throughput.png}
\caption{End-to-end redaction performance: (a) throughput and
(b) 95th-percentile latency versus request arrival rate, and
(c) PQCH adaptations and SFMR frontier cost versus request locality.}
\label{fig:exp1_redaction}
\end{figure}

Fig.~\ref{fig:exp1_redaction}(a) evaluates throughput as the offered
load increases. The baselines primarily process authorization and/or
redaction operations per request, whereas VeRedact-PQ amortizes
committee authorization across each CORA-selected set and performs
threshold PQCH adaptation once per affected transaction batch.
Accordingly, the expected advantage should become most visible near
saturation, where repeated threshold operations dominate processing
cost.

Fig.~\ref{fig:exp1_redaction}(b) evaluates the corresponding latency
trade-off. Per-Request avoids selection delay and may therefore have
lower latency under light load. CORA introduces bounded waiting up to
$T_{\max}$ to obtain useful execution sharing, but prevents indefinite
delay by making sufficiently old requests mandatory. Under increasing
load, more compatible requests become available within the waiting
window, allowing CORA to improve amortization without requiring a
fixed batch to fill. The difference between Fixed-Batch and CORA
isolates the benefit of locality-aware selection, whereas the
difference between CORA-No-SFMR and full VeRedact-PQ isolates the
execution savings obtained by SFMR.

Fig.~\ref{fig:exp1_redaction}(c) examines why these savings occur.
Independent processing requires approximately one PQCH adaptation per
redaction. In VeRedact-PQ, the number instead scales with
$|\Omega_e|$, while Merkle processing depends on the shared frontier
size $F_e$. As $s$ increases, requests are increasingly concentrated
within common transaction batches, allowing CORA to select requests
with greater batch and path overlap and SFMR to exploit that overlap.
The benefit should therefore diminish toward $s=0$ and increase with
request locality. Together, these measurements distinguish CORA's
ability to \emph{discover} execution sharing from SFMR's ability to
\emph{realize} it.

    % ---------------------------------------------------------------------

\subsubsection*{\textbf{Experiment 2: Multi-Party Authorization Performance}}

This experiment evaluates Phase~4 authorization and the benefit of
separating state-bound PQZK validation from committee approval. The
committee size is varied as
$n\in\{4,7,10,16,32\}$, with the experimental threshold configured as
$t=\lfloor2n/3\rfloor+1$, while the number of requests in a
CORA-selected set is varied as
$m\in\{1,8,32,64,128,256\}$. Authorization latency is measured from
completion of CORA selection and state revalidation to the availability
of $Auth_e^B$, and is reported both per selected set and amortized per
request.

The measured cost is decomposed into VPS-attestation verification,
state-freshness checking, authorization-root and CORA-plan construction,
and committee signing/verification. VeRedact-PQ is compared with
Schemes~\cite{ref1} and~\cite{ref34}, which provide distributed or
policy-based authorization, respectively. Schemes~\cite{ref13}
and~\cite{ref27} are excluded because they define no comparable
standalone authorization stage.

We additionally evaluate \textbf{Re-ZK}, a conservative internal
variant in which each participating committee member independently
re-verifies the PQZK proof of every selected request rather than relying
on the PQ-signed VPS attestation $\alpha_i$. This isolates the benefit
of moving expensive PQZK verification out of the committee
authorization path.

\begin{figure}[t]
\centering
\includegraphics[width=0.8\linewidth]{exp2_authorization_latency.png}
\caption{Multi-party authorization performance:
(a) amortized authorization latency per request versus selected-set
size $m$, and (b) per-set authorization latency versus committee size
$n$.}
\label{fig:exp2_auth}
\end{figure}

Fig.~\ref{fig:exp2_auth}(a) evaluates the effect of $m$. In
VeRedact-PQ, the $m$ VPS attestations are verified individually, while
the $t$ committee approvals are shared across the selected set.
Accordingly, the dominant authorization cost follows
\begin{equation}
C_{\mathrm{auth}}
=
mT_V+t(T_S+T_V)+O(m)T_H,
\end{equation}
and its amortized threshold component decreases as
$t(T_S+T_V)/m$. By contrast, schemes whose authorization work is
performed independently per request obtain little or no corresponding
batch amortization. Re-ZK additionally incurs committee-side PQZK
verification for every request, approximately proportional to
$mtT_{ZV}$ when all $t$ participating members independently verify the
proofs.

Fig.~\ref{fig:exp2_auth}(b) evaluates the committee-size trade-off.
For VeRedact-PQ, increasing $n$ also increases the configured threshold
$t$, raising committee signing and verification cost, but this cost is
incurred once per CORA-selected set and amortized across its $m$
requests. Scheme~\cite{ref1} similarly reflects its distributed
approval threshold, whereas the policy-based authorization of
Scheme~\cite{ref34} is not parameterized by committee size and therefore
serves as a non-threshold reference. The experiment thus quantifies
the cost of stronger distributed authorization and the extent to which
VeRedact-PQ amortizes that cost under larger redaction workloads.

    % ---------------------------------------------------------------------
\subsubsection*{\textbf{Experiment 3: Audit and Verification Efficiency}}

This experiment evaluates Phase~6 audit efficiency in terms of evidence
generation, communication overhead, and auditor-side verification. The
number of returned redaction records is varied as
$n_Q\in\{10,10^2,10^3,10^4\}$. Evidence sharing is controlled by
varying the number of returned records that share authorization,
transition, and checkpoint dependencies.

VeRedact-PQ is compared with Schemes~\cite{ref1} and~\cite{ref13}
where their original designs provide comparable redaction-verification
or audit functionality; unsupported operations are reported as N/A.
Two internal variants are additionally evaluated:
\begin{enumerate}
\item \textbf{Per-Record}: complete evidence is independently
transferred and verified for every returned redaction, without
dependency sharing or verification reuse; and
\item \textbf{Shared-Evidence}: common authorization, transition, and
checkpoint evidence is represented once together with a query-scoped
RAI multiproof, but shared dependencies are independently re-verified
for each dependent record.
\end{enumerate}
The full VeRedact-PQ configuration additionally applies SAVER to reuse
successfully verified shared dependencies only under matching
state-bound contexts $\chi(v)$.

We measure (i) response-generation time, (ii) audit-response size,
(iii) auditor-side verification latency, and (iv) the number of
cryptographically verified evidence nodes. Routine auditing is used
with $d_Q=0$ so that this experiment isolates evidence sharing and
verification reuse from deep PQZK escalation.

\begin{figure}[t]
\centering
\includegraphics[width=0.8\linewidth]{exp3_audit_efficiency.png}
\caption{Audit and verification efficiency:
(a) response-generation time, (b) audit-response size, and
(c) auditor-side verification latency versus the number of returned
redaction records.}
\label{fig:exp3_audit}
\end{figure}

Fig.~\ref{fig:exp3_audit}(a) compares audit-response generation as
$n_Q$ increases. Conventional per-record processing repeatedly
constructs evidence for individual redactions, whereas VeRedact-PQ
uses a query-scoped RAI multiproof and represents common authorization,
transition, and checkpoint dependencies once. Its generation cost
therefore depends on both $n_Q$ and the number of distinct shared
dependencies.

Fig.~\ref{fig:exp3_audit}(b) evaluates communication overhead.
Per-Record duplicates authentication paths and shared evidence for
each returned record, whereas Shared-Evidence and VeRedact-PQ eliminate
this duplication through $MP_{Q_j}^A$ and the verification-dependency
DAG. The reduction is expected to be greatest when many returned
records share authorization and execution contexts and to narrow as
the number of distinct dependencies approaches $n_Q$. The external
baselines provide reference points according to the audit evidence
supported by their original designs.

Fig.~\ref{fig:exp3_audit}(c) evaluates auditor-side verification.
Shared-Evidence reduces transferred evidence but still independently
re-verifies shared dependencies. VeRedact-PQ additionally verifies each
distinct shared node once under its authenticated context and reuses
the result for dependent records. For routine auditing, its cost follows
\begin{equation}
C_{\mathrm{audit}}
\approx
n_QC_{\mathrm{req}}
+a_QC_{\mathrm{auth}}
+z_QC_{\mathrm{trans}}
+c_QC_{\mathrm{cp}},
\end{equation}
where $a_Q$, $z_Q$, and $c_Q$ denote the numbers of distinct
authorization, transition, and checkpoint dependencies, respectively.
Thus, when related redactions share evidence,
$a_Q,z_Q,c_Q\ll n_Q$, reducing repeated verification without extending
a verified result across different authenticated state contexts.

Overall, the experiment distinguishes three sources of audit cost:
per-record evidence processing, communication savings from shared
evidence, and verification savings from state-bound reuse. The external
baselines quantify comparative audit performance, while the internal
variants isolate the contribution of VeRedact-PQ's evidence-sharing and
verification mechanisms.
    
    % ---------------------------------------------------------------------
 \subsubsection*{\textbf{Experiment 4: Selective Verification and Evidence Integrity}}

This experiment evaluates the cost and detection effectiveness of
Phase~6 verification under different audit depths. Unlike
Experiment~3, which evaluates evidence sharing and verification reuse,
this experiment focuses on selective PQZK escalation and the detection
of invalid audit evidence.

The number of returned records is varied as
$n_Q\in\{10,10^2,10^3,10^4\}$, while the fraction requiring independent
PQZK verification is varied as
$d_Q/n_Q\in\{0,0.01,0.05,0.10,0.25,1.0\}$.
Routine verification authenticates the VPS attestation $\alpha_i$ and
the associated request, authorization, transition, and checkpoint
dependencies. It does not require transmission or independent
verification of $\pi_i^{PQ}$. For each request selected for deep
verification, the auditor additionally retrieves the proof and checks
\begin{equation}
\mathsf{PQZK.Verify}
(vk_{\mathrm{PQZK}},x_i,\pi_i^{PQ})=1,
\qquad
H(\pi_i^{PQ})=h_i^{PQ}.
\end{equation}

Verification time is decomposed into RAI and response authentication,
request-attestation verification, committee-approval verification,
transition/PQCH verification, and PQZK verification. VeRedact-PQ is
compared with Schemes~\cite{ref1} and~\cite{ref13} for verification
operations supported by their original designs; unsupported operations
are reported as N/A.

To evaluate evidence integrity, \textbf{[1\%--10\%]} of returned
records are independently subjected to three fault classes:
\textbf{modified evidence}, where a committed field is altered;
\textbf{substituted evidence}, where evidence from another redaction or
dependency context is inserted; and \textbf{stale evidence}, where an
older authenticated state or checkpoint is replayed. We measure the
invalid-record rejection rate and valid-record acceptance rate.

\begin{figure}[t]
\centering
\includegraphics[width=0.8\linewidth]{exp4_verification_time.png}
\caption{Selective audit verification:
(a) verification time versus PQZK escalation fraction,
(b) verification-time breakdown, and
(c) invalid-record rejection rate under modified, substituted, and
stale evidence.}
\label{fig:exp4_verify}
\end{figure}

Fig.~\ref{fig:exp4_verify}(a) evaluates the cost of selective
escalation. Routine auditing corresponds to $d_Q=0$ and relies on
authenticated VPS decisions together with the required shared
authorization, transition, and checkpoint evidence. As the escalation
fraction increases, verification cost additionally includes
$d_QT_{ZV}$, approaching full independent PQZK verification when
$d_Q=n_Q$. The expected cost therefore follows
\begin{equation}
C_{\mathrm{verify}}
\approx
C_{\mathrm{routine}}+d_QT_{ZV},
\end{equation}
demonstrating the trade-off between routine verification and
request-level independent validation.

Fig.~\ref{fig:exp4_verify}(b) decomposes this cost. Routine verification
is dominated by request attestations and the distinct authorization,
transition, and checkpoint dependencies, whereas PQZK verification
becomes dominant as $d_Q/n_Q$ increases. The classical baselines provide
reference costs for the verification operations they support, while
VeRedact-PQ additionally quantifies the overhead introduced by
post-quantum signatures and selective PQZK verification.

Fig.~\ref{fig:exp4_verify}(c) evaluates verification granularity.
Modified evidence should fail its content or authenticated-path
binding, substituted evidence should fail its dependency or
state-context binding, and stale evidence should fail the applicable
version, epoch, or checkpoint checks. Because request-specific evidence
remains individually bound to the shared verification structure, an
invalid record can be rejected without invalidating unrelated records
whose dependencies remain valid. The experiment therefore evaluates
both the cost of selective verification and the ability to localize
invalid evidence without requiring full-query rejection.
    
    % ---------------------------------------------------------------------
 \subsubsection*{\textbf{Experiment 5: Permissioned Blockchain Finalization Overhead}}

This experiment evaluates the overhead introduced by PBN anchoring and
atomic redaction finalization. Because VeRedact-PQ targets a permissioned
blockchain rather than a public cryptocurrency network, we measure
transaction size and finalization latency instead of monetary gas cost.

The number of affected transaction batches is varied as
$|\Omega_e|\in\{1,2,4,8,16,32\}$. For each configuration, we measure
(i) the size of the submitted checkpoint transaction and
(ii) the time from checkpoint submission until QBFT finalization.
The experiment covers Phase~5 redaction finalization, which atomically
commits the updated batch checkpoints and authenticated SA-RLI state,
and Phase~6 RAI checkpoint anchoring. The validator count is additionally
varied as \textbf{[TBD]} to quantify consensus-related overhead.

\begin{figure}[t]
\centering
\includegraphics[width=0.8\linewidth]{exp5_blockchain_overhead.png}
\caption{Permissioned-blockchain overhead:
(a) checkpoint transaction size and
(b) finalization latency versus the number of affected transaction
batches.}
\label{fig:exp5_blockchain}
\end{figure}

Fig.~\ref{fig:exp5_blockchain}(a) evaluates the amount of state submitted
to the PBN. VeRedact-PQ anchors compact commitments and finalized
checkpoints rather than complete PBRP, SFMR, or audit evidence.
Consequently, on-chain data growth is determined primarily by the number
of affected batch checkpoints $|\Omega_e|$, while detailed supporting
evidence remains off-ledger and integrity-bound to the anchored state.

Fig.~\ref{fig:exp5_blockchain}(b) evaluates the corresponding
finalization latency. Increasing $|\Omega_e|$ enlarges the checkpoint
transaction and validator-side processing, while increasing the
validator population affects QBFT consensus overhead. This experiment
therefore quantifies the blockchain cost that remains after CORA and
SFMR reduce the number of independently finalized batch transitions.

Overall, the experiment verifies that VeRedact-PQ's blockchain layer
serves as a compact finalization and audit anchor rather than as the
primary execution environment for post-quantum cryptographic
operations.


\section{Conclusion}
\label{sec:conclusion}

This paper presented VeRedact-PQ, a scalable post-quantum framework
for policy-bound redaction and verifiable auditing in permissioned
blockchains. VeRedact-PQ combines state-bound PQZK validation,
threshold committee authorization, and distributed PQCH adaptation
with RADC and authenticated indexing. CORA identifies sharing across
validated redaction requests, while SFMR exploits this locality through
shared Merkle processing and batch-level PQCH adaptation. PBRP binds
authorization to the resulting state transition, and SAVER extends
these dependencies to state-bound audit verification with selective
PQZK escalation. The evaluation demonstrates the trade-offs among post-quantum
validation, multi-party authorization, shared redaction execution,
audit verification, and permissioned-blockchain finalization.
VeRedact-PQ achieves scalability by reducing how often expensive
operations are repeated rather than assuming cheaper post-quantum
primitives. Future work will investigate deployment-scale workloads,
adaptive committee management, stronger privacy against access-pattern
leakage, and integration with additional post-quantum redaction
constructions.
    
 \section{Acknowledgement}
    This research has been supported by the Thammasat University Research Unit in Cyber Security.
    
     \begin{thebibliography}{99}
    
    \bibitem{ref1}
    C. Li, Q. Shen and Z. Wu, "Redactable Blockchain From Decentralized Chameleon Hash Functions, Revisited," in IEEE Transactions on Computers, vol. 74, no. 6, pp. 1911-1920, June 2025, doi: 10.1109/TC.2025.3544878.
    
    
    
    
    \bibitem{ref2}
    K. Huang, X. Zhang, Y. Mu, X. Wang, G. Yang, X. Du, F. Rezaeibagha, Q. Xia, and M. Guizani,
    ``Building Redactable Consortium Blockchain for Industrial Internet-of-Things,''
    \textit{IEEE Trans. Ind. Informat.}, vol. 15, no. 6, pp. 3670--3679, Jun. 2019,
    doi: 10.1109/TII.2019.2901011.
    
    \bibitem{ref3}
    W. Wang, L. Wang, J. Duan, X. Tong, and H. Peng,
    ``Redactable Blockchain Based on Decentralized Trapdoor Verifiable Delay Functions,''
    \textit{IEEE Trans. Inf. Forensics Security}, vol. 19, pp. 7492--7507, 2024,
    doi: 10.1109/TIFS.2024.3431917.
    
    \bibitem{ref4}
    H. Guo, W. Gan, M. Zhao, C. Zhang, T. Wu, L. Zhu, and J. Xue,
    ``PriChain: Efficient Privacy-Preserving Fine-Grained Redactable Blockchains in Decentralized Settings,''
    \textit{Chinese Journal of Electronics}, vol. 34, no. 1, pp. 82--97, 2025,
    doi: 10.23919/cje.2023.00.305.
    
    \bibitem{ref5}
    D. Zhang, J. Le, X. Lei, T. Xiang, and X. Liao,
    ``Secure Redactable Blockchain With Dynamic Support,''
    \textit{IEEE Trans. Dependable Secure Comput.}, vol. 21, no. 2, pp. 717--731, 2024,
    doi: 10.1109/TDSC.2023.3261343.
    
    \bibitem{ref6}
    X. Wu, X. Du, Q. Yang, N. Wang, and W. Wang,
    ``Redactable consortium blockchain based on verifiable distributed chameleon hash functions,''
    \textit{J. Parallel Distrib. Comput.}, vol. 183, Art. no. 104777, 2024,
    doi: 10.1016/j.jpdc.2023.104777.
    
    \bibitem{ref7}
    S. Aguincha, E. Nunes, S. Eisa, and M. L. Pardal,
    ``ChainGuards: Verification of Sensed Data using Permissioned Blockchain Technology,''
    arXiv preprint arXiv:2603.20769, 2026.
    
    \bibitem{ref8}
    D. Derler, K. Samelin, D. Slamanig, and C. Striecks,
    ``Fine-Grained and Controlled Rewriting in Blockchains: Chameleon-Hashing Gone Attribute-Based,''
    in \textit{Proc. 26th Annu. Netw. Distrib. Syst. Security Symp. (NDSS)}, 2019,
    doi: 10.14722/ndss.2019.23066.
    
    \bibitem{ref9}
    G. Ateniese, B. Magri, D. Venturi, and E. R. Andrade,
    ``Redactable Blockchain -- or -- Rewriting History in Bitcoin and Friends,''
    in \textit{Proc. 2nd IEEE Eur. Symp. Security Privacy (EuroS\&P)},
    Paris, France, pp. 111--126, 2017,
    doi: 10.1109/EuroSP.2017.37.
    
    \bibitem{ref10}
    C. Peng, H. Xu, and P. Li,
    ``Redactable Blockchain Using Lattice-based Chameleon Hash Function,''
    in \textit{Proc. 2022 Int. Conf. Blockchain Technology Information Security (ICBCTIS)},
    pp. 94--98, 2022,
    doi: 10.1109/ICBCTIS55569.2022.00032.
    
    \bibitem{ref11}
    C. Wu, L. Ke, and Y. Du,
    ``Quantum resistant key-exposure free chameleon hash and applications in redactable blockchain,''
    \textit{Information Sciences}, vol. 548, pp. 438--449, 2021,
    doi: 10.1016/j.ins.2020.10.008.
    
    \bibitem{ref12}
    K. Y. Chan, L. Chen, Y. Tian, and T. H. Yuen,
    ``Reconstructing Chameleon Hash: Full Security and the Multi-Party Setting,''
    in \textit{Proc. 19th ACM Asia Conf. Computer and Communications Security (AsiaCCS)},
    pp. 1076--1091, 2024,
    doi: 10.1145/3634737.3656291.
    
    \bibitem{ref13}
    X. Zhang, Z. Cai, K. Chen, G. Ha and C. Jia, "Efficient Auditing and Querying in Verifiable Redactable Blockchain: A Lightweight VDS Protocol with Integrity Verification," 2025 IEEE 24th International Conference on Trust, Security and Privacy in Computing and Communications (TrustCom), Guiyang, China, 2025, pp. 504-511, doi: 10.1109/Trustcom66490.2025.00062.
    
    \bibitem{ref14}
    K. Huang, X. Zhang, Y. Mu, F. Rezaeibagha, and X. Du,
    ``Scalable and redactable blockchain with update and anonymity,''
    \textit{Information Sciences}, vol. 546, pp. 25--41, 2021,
    doi: 10.1016/j.ins.2020.07.016.
    
    \bibitem{ref15}
    Y. Dong, Y. Li, Y. Cheng, and D. Yu,
    ``Redactable consortium blockchain with access control: Leveraging chameleon hash and multi-authority attribute-based encryption,''
    \textit{High-Confidence Computing}, vol. 4, no. 1, Art. no. 100168, 2024,
    doi: 10.1016/j.hcc.2023.100168.
    
    \bibitem{ref16}
    Y. Zhang, Z. Ma, S. Luo, and P. Duan,
    ``Dynamic Trust-Based Redactable Blockchain Supporting Update and Traceability,''
    \textit{IEEE Trans. Inf. Forensics Security}, vol. 19, pp. 821--834, 2024,
    doi: 10.1109/TIFS.2023.3326379.
    
    \bibitem{ref17}
    X. Wang, Y. Chen, X. Zhu, C. Li, and K. Fang,
    ``A Redactable Blockchain Scheme Supporting Quantum-Resistance and Trapdoor Updates,''
    \textit{Applied Sciences}, vol. 14, no. 2, Art. no. 832, 2024,
    doi: 10.3390/app14020832.
    
    \bibitem{ref18}
    J. B. Klamti and M. A. Hasan,
    ``Revocable policy-based chameleon hash using lattices,''
    \textit{Journal of Mathematical Cryptology}, vol. 18, no. 1, Art. no. 20230012,
    pp. 152--182, 2024,
    doi: 10.1515/jmc-2023-0012.
    
    \bibitem{ref19}
    X. Huang, Y. Wang, Y. Ding, Q. Wu, C. Yang, and H. Liang,
    ``Dynamically redactable Blockchain based on decentralized Chameleon hash,''
    \textit{Digital Communications and Networks}, vol. 11, no. 3, pp. 757--767, 2025,
    doi: 10.1016/j.dcan.2024.10.013.
    
    \bibitem{ref20}
    M. Miao, X. Yang, J. Wei, G. Tian, and W. Susilo,
    ``A Verifiable and Redactable Blockchain with Lightweight Storage and Permission Supervision,''
    \textit{Information}, vol. 17, no. 2, Art. no. 176, 2026,
    doi: 10.3390/info17020176.
    
    \bibitem{ref21}
    S. Fugkeaw, S. Sungchai, S. Nakprame, and P. Sreekongpan,
    ``Enabling Secure and Scalable GDPR-Compliant Blockchain-Based e-KYC With Efficient Redaction,''
    \textit{IEEE Access}, vol. 13, pp. 136834--136853, 2025,
    doi: 10.1109/ACCESS.2025.3594656.
    
    
    \bibitem{ref22}
    C. G. Harris, "A Redactable Blockchain Architecture for Regulatory Compliance," 2025 IEEE International Conference on Blockchain (Blockchain), Zhengzhou, China, 2025, pp. 1-6, doi: 10.1109/Blockchain67634.2025.00010.
    
    \bibitem{ref23}
    Z. Wu, L. Wang, X. Zhang and X. Feng, "EMT: Extended Merkle Tree Structure for Inserted Data Redaction in Permissioned Blockchain," in IEEE Transactions on Network Science and Engineering, vol. 12, no. 4, pp. 3025-3038, July-Aug. 2025, doi: 10.1109/TNSE.2025.3555979.
    
    \bibitem{ref24}
    J. W. Heo, G. Ramachandran and R. Jurdak, "Decentralised Redactable Blockchain: A Privacy-Preserving Approach to Addressing Identity Tracing Challenges," 2024 IEEE International Conference on Blockchain and Cryptocurrency (ICBC), Dublin, Ireland, 2024, pp. 215-219, doi: 10.1109/ICBC59979.2024.10634438.
    
    
    \bibitem{ref25}
    T. Sengupta, S. Chakraborty, and S. Sural,
    ``ReAcct: Redaction control for interoperable blockchains,''
    in \textit{Proc. 7th Conf. Blockchain Research Appl. Innovative
    Networks Services (BRAINS)}, Zurich, Switzerland, 2025,
    pp. 1--10, doi: 10.1109/BRAINS67003.2025.11302947.
    
    \bibitem{ref26}
    L. Yin, J. Shi, B. Xie and Y. Liu, "Rchain: A Universal and Efficient Redactable Blockchain Scheme," 2025 IEEE International Conference on Blockchain (Blockchain), Zhengzhou, China, 2025, pp. 235-240, doi: 10.1109/Blockchain67634.2025.00038.
    
    \bibitem{ref27}
    Z. Liu, B. Zhou and Y. Zhao, "Enhancing Redactable Blockchain With Robust and Efficient Threshold Redaction," in IEEE Transactions on Dependable and Secure Computing, vol. 23, no. 2, pp. 3238-3251, March-April 2026, doi: 10.1109/TDSC.2025.3634512.
    
    \bibitem{ref28}
    F. Wang, R. Dong, J. Cui, Q. Zhang and H. Zhong, "Blockchain-Based Anonymous and Accountable Data Redaction for the Industrial Internet of Things," in IEEE Transactions on Cloud Computing, doi: 10.1109/TCC.2026.3725253.
    
    \bibitem{ref29}
    S. Biswas, S. Chakraborty, and R. S. Chakraborty,
    ``Redacting without regret: A dual-layer framework for
    policy-compliant blockchains,''
    in \textit{Proc. 7th Conf. Blockchain Research Appl. Innovative
    Networks Services (BRAINS)}, Zurich, Switzerland, 2025,
    pp. 1--9, doi: 10.1109/BRAINS67003.2025.11302930.
    
    \bibitem{ref30}
    C. Li, Q. Shen and Z. Wu, "Redactable Blockchain From Decentralized Chameleon Hash Functions, Revisited," in IEEE Transactions on Computers, vol. 74, no. 6, pp. 1911-1920, June 2025, doi: 10.1109/TC.2025.3544878.
    
    \bibitem{ref31}
    Z. Yuqi, "Policy-Hiding Chameleon Hash with Traceability for Redactable Blockchains," 2025 22nd International Computer Conference on Wavelet Active Media Technology and Information Processing (ICCWAMTIP), Chengdu, China, 2025, pp. 1-4, doi: 10.1109/ICCWAMTIP68645.2025.11352637.
    
    \bibitem{ref32}
    S. Hu, M. Li, J. Weng, J. -N. Liu, J. Weng and Z. Li, "IvyRedaction: Enabling Atomic, Consistent and Accountable Cross-Chain Rewriting," in IEEE Transactions on Dependable and Secure Computing, vol. 21, no. 4, pp. 3883-3900, July-Aug. 2024, doi: 10.1109/TDSC.2023.3339675.
    
    \bibitem{ref33}
    L. Xue, H. Huang, F. Xiao, Q. Li and W. Wang, "A Controllable, Publicly Auditable, and Redactable Blockchain With a Main–Auxiliary Architecture," in IEEE Transactions on Dependable and Secure Computing, vol. 23, no. 2, pp. 1847-1864, March-April 2026, doi: 10.1109/TDSC.2025.3620854.
    
    \bibitem{ref34}
    J. Xue et al., "Attribute-Based Policy-Hiding Redactable Blockchain With Authorizable Verification for Energy Internet," in IEEE Internet of Things Journal, vol. 12, no. 15, pp. 29570-29583, 1 Aug.1, 2025, doi: 10.1109/JIOT.2025.3569699.
    
    
    \bibitem{ref35}
    K. Huang, X. Li, F. Rezaeibagha, L. Zhang and X. Zhang, "Time Updatable Policy-Based Chameleon Hash for Traceable and Accountable Redactable Blockchain," in IEEE Transactions on Information Forensics and Security, vol. 21, pp. 1470-1483, 2026, doi: 10.1109/TIFS.2026.3655919.
    
    \bibitem{ref36}
    J. Duan, W. Wang, L. Wang and L. Gu, "Controlled Redactable Blockchain Based on T-Times Chameleon Hash and Signature," in IEEE Transactions on Information Forensics and Security, vol. 19, pp. 7560-7572, 2024, doi: 10.1109/TIFS.2024.3436925. 
    
    \bibitem{ref37}
    W. -C. Kuo, D. -R. Lin and S. -X. Chen, "Identity-Based Chameleon Hashing from the SIS Assumption in a Central Authority Model," 2025 IEEE Conference on Dependable and Secure Computing (DSC), Taipei, Taiwan, 2025, pp. 1-6, doi: 10.1109/DSC65356.2025.11260874.
    
    \bibitem{ref38}
    M. Shirmohammadi, A. Islam and H. Karimipour, "An Intent-Based Networking Framework for Secure and Privacy-Compliant Machine Unlearning Using Meta-Learning and Redactable Blockchain," in IEEE Internet of Things Journal, vol. 13, no. 9, pp. 18167-18181, 1 May1, 2026, doi: 10.1109/JIOT.2025.3638966. 
    \end{thebibliography}
    
    \end{document}